\section{Conventions and foundations}\label{fd:sec}

\paragraph{Purpose of this document.}
This document collects human-readable proofs of results of a research project on the
Erd\H{o}s--Szekeres problem whose formal versions have been checked with the interactive theorem prover
Lean~4 and its mathematical library Mathlib. The Lean checks remain the primary evidence for these results. The
written proofs are meant to make the arguments accessible to a mathematician who does not read Lean:
every statement that corresponds to a Lean declaration is stated exactly as in Lean, with the name of
the declaration and its file and line, and it receives a complete proof that can be checked without
Lean. There are two exceptions. A statement labelled computer-assisted is reduced to a finite
computation, which is stated precisely and need not be carried out by hand; where this document also
gives a complete written proof by hand, the label is qualified as ``computer-assisted in Lean'' (see the
rigor labels below). A few
published theorems that the Lean code uses (for example Baek's theorem, Theorem~\ref{mb:thm:baek}) are
cited in the exact form in which they are used and are not reproved. A convenient variant of a Lean
statement, for example one in a different generality, is labelled as a variant. These written proofs
have not been refereed. Where a Lean theorem assumes propositions that are not proved in Lean (for
example verdicts of SAT solvers), the statement says so and the theorem is presented as a conditional
statement.

\paragraph{Files, commit and versions.}
All files belong to the project repository at commit \dirn{7e32ee259} (branch
\dirn{es7-night-2026-11-01}; Lean package \dirn{erdos_szekeres}). A path that begins with
\dirn{ErdosSzekeres/}, \dirn{docs/}, \dirn{experiments/} or \dirn{scripts/} is relative to the
repository root, and so are the files \dirn{lean-toolchain}, \dirn{lake-manifest.json} and
\dirn{lakefile.toml} at the top of the repository. The directory \dirn{ErdosSzekeres/} is the \emph{main
library}, the default build target of the package. Every other path, for example
\dirn{LAYERLEAN/LayerLemma.lean:298}, is relative to the research-records directory, whose path from the
repository root is \dirn{cloud/es7-2026-10-19/results}. The Lean files of the research records are not
part of the build target; they are checked separately, for example with \texttt{lake env lean} run from
the repository root (as \dirn{LAYERLEAN/logs/axioms_recheck_20261106/README.txt} records for a copy of
\dirn{LAYERLEAN/LayerLemma.lean} followed by \texttt{\#print axioms} commands), and each section names the
build and axioms logs it relies on. After a section has cited a file with its full path, it may refer to
it by its bare file name. A suffix \dirn{:N} is a line number; for a declaration it is the line that
carries the declaration keyword (\lean{theorem}, \lean{lemma}, \lean{def}, \lean{abbrev},
\lean{structure} or \lean{inductive}, possibly preceded by an attribute such as \texttt{@[simp]} or a
modifier such as \lean{noncomputable} or \lean{private}). A bare \dirn{:N} written after a file has been
cited refers to line~$N$ of that file. Lean declarations
are given by their full names in statement heads, for
example \lean{ErdosSzekeres.Suk.convexPosition_of_four_subsets}; in running text we sometimes omit the
common prefix \lean{ErdosSzekeres.} (so \lean{Suk.orient_reflectY} is
\lean{ErdosSzekeres.Suk.orient_reflectY}). Names in a namespace outside \lean{ErdosSzekeres}, such as
\lean{MirrorOdd.hasConvexNGon_of_le} in the top-level namespace \lean{MirrorOdd}, are always written in
full. Mathlib declarations are named as in Mathlib. Lean expressions are quoted in typewriter type, for
example \texttt{Finset Point}.
The Lean files were checked with Lean 4.24.0 (file \dirn{lean-toolchain}) and Mathlib at revision
\dirn{f897ebcf} (full hash \dirn{f897ebcf72cd16f89ab4577d0c826cd14afaafc7}, file
\dirn{lake-manifest.json}).

\paragraph{Rigor labels.}
We use exactly three labels; the first and the third also occur in a qualified form, defined below.
\begin{itemize}
\item \emph{Lean-checked}: the statement is a Lean declaration that Lean~4.24.0 with Mathlib
  \dirn{f897ebcf} accepted, using no axioms other than the three standard axioms \lean{propext},
  \lean{Classical.choice} and \lean{Quot.sound}; for a declaration whose proof uses Lean's
  \texttt{native\_decide}, directly or through a declaration on which it depends, the two further axioms
  \lean{Lean.ofReduceBool} and \lean{Lean.trustCompiler} that \texttt{native\_decide} adds are allowed as
  well, and the section that cites such a declaration names them (no declaration cited in this section is
  of this kind). The evidence is the saved output of
  Lean's command \texttt{\#print axioms}, for the declaration itself or for a declaration whose statement
  or proof uses it (this command reports every axiom in the dependency closure, which contains the
  declarations used in statements as well as in proofs). Each section names the build and axioms logs
  on which this label rests. Where a declaration is not covered by such an output, the section says so
  and names the evidence used instead, which is of one of two kinds:
  \begin{enumerate}[label=(\alph*)]
  \item evidence that the current version of the file that contains the declaration compiled without
    errors (exit status~$0$), namely a saved compile or build log of that file, or a project record of
    such a build when saved Lean output of a build that includes the file (for example an axiom audit)
    exists as well; together with a text search of the file and of the project files it imports that
    finds no \texttt{sorry} (a declaration that uses \texttt{sorry} compiles with a warning only, so the
    search, not the exit status, excludes it), no \texttt{admit}, no \texttt{axiom} declaration, no
    \texttt{native\_decide}, no \texttt{+native} (as in \texttt{decide +native}, of which
    \texttt{native\_decide} is a synonym), no \lean{Lean.ofReduceBool} or \lean{Lean.ofReduceNat} (the
    axioms through which values computed by Lean's compiler or interpreter enter a proof) and no
    \lean{debug.skipKernelTC} (an option that skips the kernel's type check);
  \item a \texttt{\#guard\_msgs} command, compiled in a build whose log is saved, that pins the output of
    \texttt{\#print axioms} for the declaration or for a declaration whose statement or proof uses it.
  \end{enumerate}
  A text search of every Lean file of the repository that this document cites, and of the project files
  these import, finds none of the strings \texttt{+native}, \texttt{ofReduceBool}, \texttt{ofReduceNat}
  and \texttt{skipKernelTC}, so every application of~(a) in this document remains valid, also where it
  reports only a search for \texttt{sorry}, \texttt{admit}, \texttt{axiom} declarations and
  \texttt{native\_decide}.
  If no Lean output of any kind (compile or build log, axiom output, pinned message) is saved for a
  build that includes a file, and the label of its declarations rests only on the statement of a project
  record such as a README, we write \emph{Lean-checked according to the record}.
\item \emph{Written proof (not refereed)}: every proof in this document. It is meant to be checkable by a
  mathematician without Lean; it has not been refereed.
\item \emph{Computer-assisted}: a statement that rests on a finite computation carried out by a machine.
  This is the case, first, when the Lean proof of the statement, or of one of its parts, is itself a
  finite computation: that statement or part is decided by kernel evaluation of a decision procedure (Lean's
  \texttt{decide} or \texttt{decide +kernel}, which add no axiom) or by compiled evaluation (Lean's
  \texttt{native\_decide}, which adds the axioms \lean{Lean.ofReduceBool} and \lean{Lean.trustCompiler}
  and trusts the compiler), or it is obtained by evaluating generated data or by replaying a certificate
  produced by an external program (for example an LRAT refutation found by a SAT solver, replayed in Lean
  and checked by the kernel). It is the case, second, when the written proof of the statement uses a
  machine computation that it does not carry out by hand, for example by applying a statement of the
  first kind that has no written proof by hand. For a computer-assisted statement we give the
  mathematical reduction to the computation and describe the computation precisely; we do not claim to
  have carried out by hand a computation for which no written proof by hand is given. If the Lean proof
  of a statement is a finite computation and this document also gives a complete written proof of the
  statement that does not use the computation, the statement is labelled \emph{computer-assisted in
  Lean}, usually with a pointer to the written proof (as in ``computer-assisted in Lean; written proof by
  hand''); that written proof can be checked without a machine. A Lean proof that uses \texttt{decide}
  only for side facts inside the proof of a statement that is not itself a finite computation (such as
  $0<1$ in $\mathrm{Fin}\,4$, or closed instances that the written proof checks by hand) does not make the
  statement computer-assisted. Certificates may also be checked by programs outside Lean. For the SAT
  certificates of the project only the checker \texttt{cake\_lpr} counts, and a verdict of it counts only
  if it prints the line \texttt{s VERIFIED UNSAT}; verdicts of \texttt{lrat-check} alone are not counted
  (\dirn{LRATCHECK20261105/README.md}; \dirn{WRITEUP/SYMMETRIC.md}, section~3). A section may report that
  a result of the project's records, such as an unsatisfiability verdict, is computer-assisted; this is a
  claim of the records, which this document does not prove, and a Lean theorem that takes such a result
  as a hypothesis is presented as conditional on it.
\end{itemize}
A statement, or a part of a statement, whose head names a declaration of the project in the form
``Lean: name, file:line'' is Lean-checked in this sense. For a definition this records only that Lean
accepted the definition: a definition asserts nothing, and what this document claims for it is that the
written definition restates the Lean definition. A part marked ``written step'' does not carry this
label, even where the text names a Lean counterpart of it elsewhere, and a Mathlib declaration named in a
head only identifies the Mathlib counterpart of a statement that is proved here. A Lean theorem that
takes unproved propositions as hypotheses is Lean-checked only as the conditional statement it is, and
the hypotheses are named. We use no other words of this kind (such as ``verified'' or ``certified'') for
the results of this document.

\paragraph{Lean's arithmetic, written explicitly.}
Lean statements often use natural numbers inside real expressions and subtract natural numbers. We make
both explicit. A natural number $m$ used as a real number is written $(m:\R)$ whenever the cast matters;
for example the Lean angle \texttt{2 * Real.pi / m} with \lean{m} a natural number is $2\pi/(m:\R)$. Real
division in Lean is total, with $x/0=0$; hence $2\pi/(m:\R)=0$ when $m=0$. Subtraction of natural
numbers in Lean is truncated ($a-b=0$ when $a<b$); whenever a Lean statement subtracts natural numbers we
write the difference as $\max(a-b,0)$, or as $a\mathbin{\dot-}b=\max(a-b,0)$ where a section defines this
symbol, or we state the inequality $b\le a$ that makes it an ordinary difference.
$\mathrm{Fin}\,n=\{0,1,\dots,n-1\}$ carries its natural order. For a map $f$ we write $f^{[n]}$ for its
$n$-th iterate ($f^{[0]}=\mathrm{id}$, $f^{[n+1]}=f^{[n]}\circ f$), as Lean does; the recursion
$f^{[n+1]}=f\circ f^{[n]}$ defines the same maps (see the proof of Lemma~\ref{fd:lem:inv-image}(c)).
Mathlib's \lean{Function.minimalPeriod} $f$ $x$ is the least $n>0$ with $f^{[n]}(x)=x$, or $0$ if there
is no such $n$.

\paragraph{Default notation.}
Points, finite sets and $\card S$ are as in Definition~\ref{fd:def:point}; $\orient$ is
Definition~\ref{fd:def:orient}; $\conv$ is the convex hull of Definition~\ref{fd:def:conv}. For a finite
set $X$ of points, $\hv(X)=\{x\in X: x\notin\conv(X\setminus\{x\})\}$ is the set of hull vertices
(Lean \lean{hullVerts}) and $\peel(X,0)=X$, $\peel(X,n+1)=\peel(X,n)\setminus\hv(\peel(X,n))$
(Lean \lean{peel}); these two notions are defined formally, with their Lean names, in
Definitions~\ref{ll:def:hv} and~\ref{ll:def:peel}. ``$S$ is $K$-free'' means
$\neg\,\mathrm{HasConvexNGon}\,K\,S$ (Definition~\ref{fd:def:ngon}). We write $\rho_{c,\theta}$ for the
rotation $\mathrm{rotAbout}\,c\,\theta$ of Definition~\ref{fd:def:rot}; rotations act counterclockwise
(Lemma~\ref{fd:lem:rot}(e)). In this section $R_\theta$ always denotes the rotation matrix of
Lemma~\ref{fd:lem:rot}, not a map. The notation of this paragraph and of the previous one is the
default for the whole document. A later section may write the same objects differently, for example
$\mathrm{rotAbout}(c,\theta)$ or $\mathrm{rotAbout}_{c,\theta}$ for $\rho_{c,\theta}$, a letter such as
$R$ with subscripts for the rotation map itself and another letter for its matrix, $f^n$ for $f^{[n]}$,
or $\mathrm{HasConvexNGon}$ in another typeface; such a section defines its symbols where it introduces
them.

\paragraph{This section.}
We define points, the orientation function, general position, convex position, convex $n$-gons,
$K$-freeness, forcing, invariance and rotations exactly as the Lean code does, and we prove the basic
facts that all later sections use: the symmetries of $\orient$ and its behaviour under affine maps; the
equivalence of $\orient(a,b,c)=0$ with collinearity; the hull form of convex position; monotonicity;
invariance of convex and general position under injective affine maps; Carath\'eodory's theorem in the
plane; the four-point criteria for convex position; the criterion that points whose increasing triples
are all counterclockwise are in convex position; and the separation of a point from a finite set. This
section relies on no other section. Its sources, with the first 12 hexadecimal digits of their SHA-256
hashes, are listed in the following table, the table of sources.
\begin{longtable}{ll}
\toprule
file & SHA-256 (prefix)\\
\midrule
\endfirsthead
\toprule
file & SHA-256 (prefix)\\
\midrule
\endhead
\bottomrule
\endfoot
\dirn{ErdosSzekeres/Definitions.lean} & \texttt{2425500eefd1}\\
\dirn{ErdosSzekeres/Orientation.lean} & \texttt{f67955ac7091}\\
\dirn{ErdosSzekeres/OrientationGeometry.lean} & \texttt{c7db9bbbb65e}\\
\dirn{ErdosSzekeres/Basic.lean} & \texttt{13667aa2d3dc}\\
\dirn{ErdosSzekeres/Suk/FourPoint.lean} & \texttt{f3783fe9766f}\\
\dirn{ErdosSzekeres/Suk/SupportingLines.lean} & \texttt{6a8e12677685}\\
\dirn{ErdosSzekeres/Suk/OrientationClique.lean} & \texttt{cf63d5219b47}\\
\dirn{ErdosSzekeres/ConvexCycleCuts.lean} & \texttt{b4ad0ba9e583}\\
\dirn{ErdosSzekeres/Suk/Reflection.lean} & \texttt{8b99cbf5fa2b}\\
\dirn{LAYERLEAN/LayerLemma.lean} & \texttt{afb8542dceb0}\\
\dirn{SYMLEAN/SymmetricReduction.lean} & \texttt{6d5ec90a4215}\\
\dirn{SYMLEAN/AffineReduction.lean} & \texttt{41cfc0b60aec}\\
\dirn{LEANEXT/orbit/OrbitConsistency.lean} & \texttt{9b843180ac38}\\
\dirn{MIRRORETV/lean/MirrorMain.lean} & \texttt{ca623e55a1b9}\\
\dirn{SYMLEAN/mirror/ReductionThreeClasses.lean} & \texttt{0cc83994fe8c}\\
\end{longtable}

\paragraph{Axioms evidence.}
{\raggedright
In this paragraph, ``line~$N$'' without a file name refers to the log, audit or report named last
before it.
A declaration said to be used ``at'' a file and line is used on that line, inside the proof named; the
line is not the line of its own declaration.
\begin{itemize}
\item \emph{Main library, printed theorems.} The saved axiom audit
  \dirn{docs/validation/axiom-audit-lower-bound-2026-10-11.txt} reports exactly the axioms \lean{propext},
  \lean{Classical.choice}, \lean{Quot.sound} for \lean{ErdosSzekeres.hasConvexNGon_of_positive_cycle}
  (line~87), \lean{ErdosSzekeres.Suk.convexPosition_of_four_subsets} (line~103),
  \lean{ErdosSzekeres.Suk.convexPosition_of_strict_supports} (line~113),
  \lean{ErdosSzekeres.Suk.convexPosition_four_of_positive} (line~114),
  \lean{ErdosSzekeres.Suk.convexPosition_of_positive_triples} (line~119),
  \lean{ErdosSzekeres.HasConvexNGon.linearEquiv} (line~175) and
  \lean{ErdosSzekeres.Suk.generalPosition_reflectY} (line~176).
\item \emph{Main library, declarations used by these theorems.} Since \texttt{\#print axioms} reports every
  axiom used anywhere in the dependency closure of a declaration, these lines also cover the declarations
  used in those proofs: \lean{Suk.convexPosition_of_subsets_card_le_four}, \lean{convexPosition_iff},
  \lean{ConvexPosition.mono}, \lean{Suk.convex_orient_positive}, \lean{Suk.orient_affine_combination},
  \lean{Suk.orient_support_sum}, \lean{Suk.reflectY}, \lean{Suk.orient_reflectY},
  \lean{generalPosition_iff_orient_ne_zero} and \lean{orient_eq_zero_iff_collinear}
  (\lean{Suk.generalPosition_reflectY} uses \lean{generalPosition_iff_orient_ne_zero} at
  \dirn{ErdosSzekeres/Suk/Reflection.lean:55}, and the proof of the latter uses
  \lean{orient_eq_zero_iff_collinear}).
\item \emph{Main library, further theorems of the audit.} Other theorems printed in the same audit, with
  the same three axioms, cover further declarations through their statements or proofs:
  \begin{itemize}
  \item line~725, \lean{conjecturedBound_le_of_forces}, uses \lean{ForcesConvexNGon.mono} at
    \dirn{ErdosSzekeres/LowerBound/Construction.lean:118};
  \item line~732, \lean{esNumber_three} (declared at \dirn{ErdosSzekeres/LowerBound/Number.lean:40}), uses
    \lean{forces_triangle}, whose proof uses \lean{forcesConvexNGon_iff_exact} and
    \lean{convexPosition_of_card_three} at \dirn{ErdosSzekeres/Basic.lean:73} and \dirn{:75}; the proof
    of \lean{forcesConvexNGon_iff_exact} uses \lean{GeneralPosition.mono} and \lean{HasConvexNGon.mono} at
    \dirn{ErdosSzekeres/Basic.lean:45};
  \item line~720, \lean{LowerBound.exists_cupCapFree}, uses \lean{LowerBound.generalPosition_translate}
    at \dirn{ErdosSzekeres/LowerBound/Construction.lean:54}, whose proof uses \lean{orient_translate} at
    \dirn{ErdosSzekeres/LowerBound/Placement.lean:20};
  \item line~624, \lean{Omission.Configuration.orient_ne_of_distinct}, uses \lean{orient_swap} and
    \lean{orient_cycle} at \dirn{ErdosSzekeres/OrientationObstructions.lean:64} and \dirn{:68};
  \item line~132, \lean{Suk.AdjacentCapFrames.glue_chain_caps}, uses \lean{ConvexPosition.hasConvexNGon}
    at \dirn{ErdosSzekeres/Suk/SpikeChain.lean:140};
  \item line~49, \lean{VisibleEdgeInsertion.convexPosition_insert_of_one_visible_edge}, uses
    \lean{Suk.convex_orient_nonnegative} at \dirn{ErdosSzekeres/VisibleEdgeInsertion.lean:77};
  \item line~178, \lean{Suk.horizontal_finite_density_bound}, uses \lean{Suk.hasConvexNGon_of_reflectY} at
    \dirn{ErdosSzekeres/Suk/FiniteDensityBound.lean:66};
  \item line~88 prints \lean{convexCycleCut_of_avoids}, whose statement contains \lean{ConvexCycleCut};
  \item line~529 prints \lean{Omission.BoundaryGeometry.exposed_not_in_hull} (declared at
    \dirn{ErdosSzekeres/Omission/BoundaryGeometry.lean:30}), which the Coverage paragraph names.
  \end{itemize}
  The definitions \lean{Point}, \lean{orient}, \lean{GeneralPosition}, \lean{ConvexPosition},
  \lean{HasConvexNGon} and \lean{ForcesConvexNGon} occur in the statements of the printed theorems above
  (for example in those of lines 87, 103, 176 and 725), so they are covered as well. Every main-library
  file named in this paragraph is unchanged since commit \dirn{7f0868fdb} (11~October 2026), which added
  the audit.
\item \emph{Main library, remaining declarations.} The remaining main-library declarations named in this
  section, \lean{HasConvexNGon.card_le}, \lean{Suk.convexPosition_of_negative_triples} and the three
  simplification lemmas \lean{Suk.reflectY_x}, \lean{Suk.reflectY_y}, \lean{Suk.reflectY_reflectY} of
  Remark~\ref{fd:rem:affine-cp-lean}, are not covered by a printed line that we identified. For them the
  evidence is of kind~(a) of the rigor labels: the build record \dirn{docs/VALIDATION.md:73} (on
  11~October 2026 the main library, which contains their files, builds with \texttt{lake build
  -{}-wfail}, which fails on any warning, in particular on the warning that Lean issues for a declaration
  that uses \texttt{sorry}); the audit above, saved Lean output of the same library, whose printed
  theorems include theorems of \dirn{ErdosSzekeres/Suk/OrientationClique.lean} and
  \dirn{ErdosSzekeres/Suk/Reflection.lean} and depend on declarations of \dirn{ErdosSzekeres/Basic.lean};
  and a text search of their files and of the main-library files these import
  (\dirn{ErdosSzekeres/Suk/BaseOrder.lean}, \dirn{ErdosSzekeres/Suk/BaseChart.lean},
  \dirn{ErdosSzekeres/Suk/TriangleOrder.lean}, \dirn{ErdosSzekeres/Suk/CupsCaps.lean} and the four files
  of \dirn{ErdosSzekeres/Vendor/ErdosSzekeresConvex/}, besides the files in the table of sources), which
  finds no \texttt{sorry}, \texttt{admit}, \texttt{axiom} declaration or \texttt{native\_decide}, and no
  \texttt{implemented\_by} or \texttt{extern} attribute.
\item \emph{The layer-lemma file.} The review report \dirn{LEANEXT/review/REPORT.md} (item~1) records that
  \dirn{LAYERLEAN/LayerLemma.lean}, in its current version (1148 lines, SHA-256 prefix \texttt{afb8542d}),
  compiles with exit code~0 and no messages. The re-check
  \dirn{LAYERLEAN/logs/axioms_recheck_20261106/out.txt} records exit status~0 for a file whose first 1148
  lines are byte-identical to \dirn{LAYERLEAN/LayerLemma.lean}, followed by 19 \texttt{\#print axioms}
  commands; the only Lean output in it is the 19 printed axiom lines. (The log
  \dirn{LAYERLEAN/logs/build.log} names no file hash; it was committed with the first, 454-line version of
  the file and has not been updated since, and we do not use it.) The review log
  \dirn{LEANEXT/review/axioms.log}, which prints the axioms of the theorems of
  \dirn{LAYERLEAN/LayerLemma.lean}, lists \lean{ErdosSzekeres.LayerLemma.iterate_mem} with \lean{propext}
  and \lean{Quot.sound} only (line~1) and \lean{ErdosSzekeres.LayerLemma.rotAbout_eq_crot} (line~22),
  \lean{ErdosSzekeres.LayerLemma.rotAbout_injective} (line~23),
  \lean{ErdosSzekeres.LayerLemma.hullVerts_eq_orbitFS_rotAbout} (line~26; its proof uses
  \lean{ErdosSzekeres.LayerLemma.rotAff} at \dirn{LAYERLEAN/LayerLemma.lean:368}) and
  \lean{ErdosSzekeres.LayerLemma.rotAbout_iterate} (line~33) with the three standard axioms.
\item \emph{The symmetric and affine reduction files.} \dirn{SYMLEAN/logs/runs.tsv} records that
  \dirn{SYMLEAN/SymmetricReduction.lean} (SHA-256 prefix \texttt{6d5ec90a}, the current file) compiles
  with exit code~0 and no output (line~2), and so does \dirn{SYMLEAN/AffineReduction.lean} (prefix
  \texttt{41cfc0b6}, the current file; line~4). The file \dirn{SYMLEAN/AffineReductionAxioms.lean} is
  \dirn{SYMLEAN/AffineReduction.lean} followed by nine \texttt{\#print axioms} commands and one
  \texttt{\#check} command; it compiles with exit code~0 (\dirn{SYMLEAN/logs/runs.tsv}, line~5), and its
  output \dirn{SYMLEAN/logs/AffineReductionAxioms.out} lists with the three standard axioms
  \begin{itemize}
  \item \lean{ErdosSzekeres.SymReduction.heptagon_of_nontrivial_affine_coords} (lines~9--11), whose proof
    uses \lean{ErdosSzekeres.SymReduction.hasConvexNGon_of_affine_coords_symmetry} at
    \dirn{SYMLEAN/AffineReduction.lean:1528}, whose proof in turn uses
    \lean{ErdosSzekeres.SymReduction.injective_of_affine} at \dirn{SYMLEAN/AffineReduction.lean:1462};
  \item \lean{ErdosSzekeres.SymReduction.generalPosition_image_affineEquiv} (lines~13--15), whose proof
    uses \lean{ErdosSzekeres.SymReduction.generalPosition_image_affine} at
    \dirn{SYMLEAN/AffineReduction.lean:1295}, whose proof in turn uses
    \lean{ErdosSzekeres.SymReduction.orient_affine} at \dirn{SYMLEAN/AffineReduction.lean:1281};
  \item \lean{ErdosSzekeres.SymReduction.hasConvexNGon_image_affine} (line~16), whose proof uses
    \lean{ErdosSzekeres.SymReduction.convexPosition_image_affine} at
    \dirn{SYMLEAN/AffineReduction.lean:1303}.
  \end{itemize}
\item \emph{The orbit-consistency file.} \dirn{LEANEXT/orbit/logs/build.log} records exit code~0 for the
  build of \dirn{LEANEXT/orbit/OrbitConsistency.lean}; the log contains only the exit code, and the report
  \dirn{LEANEXT/orbit/REPORT.md} that accompanies these logs identifies the checked file (line~3: 429
  lines, SHA-256 prefix \texttt{9b843180}, the current file). \dirn{LEANEXT/orbit/logs/axioms.log} lists
  the two theorems of that file named in this section,
  \lean{ErdosSzekeres.OrbitConsistency.orient_rotAbout_inv} (line~1) and
  \lean{ErdosSzekeres.OrbitConsistency.orient3_perm} (line~5), with the three standard axioms. From
  \dirn{SYMLEAN/SymmetricReduction.lean} and \dirn{LEANEXT/orbit/OrbitConsistency.lean} we otherwise cite
  only definitions.
\item \emph{The mirror files.} \lean{MirrorOdd.hasConvexNGon_of_le} is used at
  \dirn{MIRRORETV/lean/MirrorMain.lean:2520}, in the proof of \lean{MirrorOdd.mirror_odd_general}
  (declared at \dirn{MIRRORETV/lean/MirrorMain.lean:2498}), which the compile log
  \dirn{MIRRORETV/lean/MirrorMain.out} of that file (warnings, no errors) lists with the three standard
  axioms (line~858). Its verbatim copy in \dirn{SYMLEAN/mirror/ReductionThreeClasses.lean} is used in the
  same way, at \dirn{SYMLEAN/mirror/ReductionThreeClasses.lean:4079}, in the proof of the copy of
  \lean{MirrorOdd.mirror_odd_general} (declared at \dirn{SYMLEAN/mirror/ReductionThreeClasses.lean:4057}),
  which line~11 of \dirn{SYMLEAN/mirror/axioms.out} lists with the three standard axioms; and
  \dirn{SYMLEAN/mirror/README.md} records that this file compiles with exit code~0 and no errors.
\item \emph{Text search.} A text search of the six research files of the table of sources finds no
  \texttt{sorry}, \texttt{admit}, \texttt{axiom} declaration or \texttt{native\_decide}, and no
  \texttt{implemented\_by} or \texttt{extern} attribute.
\end{itemize}
\par}

%======================================================================
\subsection{Definitions}\label{fd:sec-defs}

\begin{definition}[Points and finite point sets; \leanref{ErdosSzekeres.Point}{ErdosSzekeres/Definitions.lean:9}]\label{fd:def:point}\sloppy
A \emph{point} is an element of $\R^2$. In Lean, \lean{Point} abbreviates the type of functions from
$\mathrm{Fin}\,2=\{0,1\}$ to $\R$; we write $p=(p_0,p_1)$, where $p_0=p(0)$ is the first
($x$-)coordinate and $p_1=p(1)$ the second ($y$-)coordinate. Addition, subtraction and multiplication by
real numbers are coordinatewise.

A \emph{finite set of points} is a Lean \texttt{Finset Point}: a finite set, without multiplicities. We use
ordinary set notation for finite sets: $S\subseteq T$, $S\cup T$, $S\setminus T$, $S\setminus\{c\}$
(Lean \texttt{S.erase c}), $\card S$ (Lean \lean{S.card}), and $f(S)=\{f(p):p\in S\}$ for a map $f$ (Lean
\texttt{S.image f}, again a finite set, with $\card{f(S)}=\card S$ when $f$ is injective). Where Lean needs
a set rather than a finite set, for example as the argument of a convex hull, it uses the coercion of $S$
to a \texttt{Set Point}, which has the same elements; we do not distinguish the two. Brace notation
$\{a,b,c\}$ denotes the set with these elements; it has fewer than three elements when two of $a,b,c$
coincide.
\end{definition}

\begin{remark}[Norm and topology]\label{fd:rem:norm}
Mathlib gives the Lean type of points its default sup norm $\max(\abs{p_0},\abs{p_1})$. No statement of
this section uses a norm or a metric; Euclidean quantities are always written out in coordinates (as in
Lemma~\ref{fd:lem:separation}). The only topological argument of this section is the compactness
argument in the proof of Lemma~\ref{fd:lem:separation}, which takes place in $\R^m$ with its usual
topology. (On $\R^2$ the sup norm and the Euclidean norm define the same, usual, topology.)
\end{remark}

\begin{definition}[Orientation; \leanref{ErdosSzekeres.orient}{ErdosSzekeres/Orientation.lean:9}]\label{fd:def:orient}\sloppy
Let $R$ be a commutative ring and let $a,b,c\in R^2$ (in Lean: functions from $\mathrm{Fin}\,2$ to
$R$). Then
\[
  \orient(a,b,c)=(b_0-a_0)(c_1-a_1)-(b_1-a_1)(c_0-a_0)\in R .
\]
Unless stated otherwise $R=\R$; Lean also evaluates $\orient$ over $\mathbb{Q}$ and $\Z$ in exact
certificate checks. Equivalently, $\orient(a,b,c)=\det\bigl(b-a\;\big|\;c-a\bigr)$ is the determinant of
the $2\times2$ matrix with columns $b-a$ and $c-a$, that is, twice the signed area of the triangle $abc$.
Positive orientation means counterclockwise; for example $\orient((0,0),(1,0),(0,1))=1>0$. For $a\ne b$
we say that $c$ lies (strictly) to the \emph{left} of the line through $a$ and $b$ directed from $a$ to
$b$ if $\orient(a,b,c)>0$, and to the \emph{right} if $\orient(a,b,c)<0$.
\end{definition}

\begin{definition}[Convex sets, convex hulls, affine maps; Mathlib \lean{Convex}, \lean{convexHull}, \lean{AffineIndependent}, \lean{AffineMap}, \lean{AffineEquiv}]\label{fd:def:conv}
A set $K\subseteq\R^2$ is \emph{convex} if $\lambda x+\mu y\in K$ for all $x,y\in K$ and all reals
$\lambda,\mu\ge0$ with $\lambda+\mu=1$ (Mathlib \lean{Convex}). The \emph{convex hull} $\conv Y$ of a set
$Y\subseteq\R^2$ is the intersection of all convex sets that contain $Y$ (Mathlib \lean{convexHull},
defined as the smallest convex superset). A \emph{convex combination} of points $y_1,\dots,y_k$
($k\ge1$) is a point $\sum_{i=1}^k\lambda_iy_i$ with all $\lambda_i\ge0$ and $\sum_i\lambda_i=1$. A finite
set $T\subseteq\R^2$ is \emph{affinely independent} if the only family of reals $(\mu_y)_{y\in T}$ with
$\sum_{y\in T}\mu_y=0$ and $\sum_{y\in T}\mu_y\,y=0$ is the zero family (equivalent to Mathlib's
\lean{AffineIndependent} for the inclusion map of $T$). An \emph{affine map} of $\R^2$ is a map
$g(x)=Ax+v$, where $A$ is a real $2\times2$ matrix (the \emph{linear part} of $g$) and $v\in\R^2$.
In Mathlib, an affine map $f$ from the plane to itself (\lean{AffineMap}) is a map $f\colon\R^2\to\R^2$
together with a linear map $f_{\mathrm{lin}}\colon\R^2\to\R^2$ (Lean \lean{f.linear}) such that
$f(w+p)=f_{\mathrm{lin}}(w)+f(p)$ for all $p,w\in\R^2$. These are exactly the affine maps just defined:
taking $p=0$ gives $f(x)=A_fx+f(0)$, where $A_f$ is the matrix of $f_{\mathrm{lin}}$; conversely,
$g(x)=Ax+v$ satisfies $g(w+p)=Aw+g(p)$, so it is a Mathlib affine map whose linear map is $x\mapsto Ax$.
A Mathlib \emph{affine equivalence} of the plane (\lean{AffineEquiv}) is a bijection $f$ of $\R^2$
together with a linear bijection $f_{\mathrm{lin}}$ of $\R^2$ with the same property; in particular it
is an injective affine map.
\end{definition}

\begin{definition}[General position; \leanref{ErdosSzekeres.GeneralPosition}{ErdosSzekeres/Definitions.lean:12}]\label{fd:def:gp}\sloppy
A set $X\subseteq\R^2$ is \emph{collinear} (Mathlib \lean{Collinear}) if the linear span of the set of
differences $\{p-q:p,q\in X\}$ (Mathlib \lean{vectorSpan}) has dimension at most~$1$ (in Lean:
\lean{Module.rank} of this span is at most $1$). A finite set $S$ of points is \emph{in general
position} if for all $a,b,c\in S$ with $a\ne b$, $a\ne c$ and $b\ne c$ the set $\{a,b,c\}$ is not
collinear.
\end{definition}

\begin{definition}[Convex position; \leanref{ErdosSzekeres.ConvexPosition}{ErdosSzekeres/Definitions.lean:17}]\label{fd:def:cp}\sloppy
A finite set $S$ of points is \emph{in convex position} if the inclusion map of $S$ into $\R^2$ is
convex independent over $\R$ (Mathlib \lean{ConvexIndependent}). Mathlib calls a family
$p\colon\iota\to\R^2$ convex independent if for every set $t\subseteq\iota$ and every $x\in\iota$,
$p(x)\in\conv p(t)$ implies $x\in t$. For the inclusion map of $S$ the index set is $S$ itself (in Lean,
the type of elements of $S$) and $p$ is the identity on it, so the definition says: \emph{for every
$T\subseteq S$ and every $x\in S$, if $x\in\conv T$ then $x\in T$.} Lemma~\ref{fd:lem:cp-hull} gives the
familiar form: no point of $S$ lies in the convex hull of the other points of $S$.
\end{definition}

\begin{definition}[Convex $n$-gons and $K$-free sets; \leanref{ErdosSzekeres.HasConvexNGon}{ErdosSzekeres/Definitions.lean:21}]\label{fd:def:ngon}\sloppy
For $n\in\N$ and a finite set $S$ of points, $\mathrm{HasConvexNGon}\,n\,S$ holds if there is a finite
set $T\subseteq S$ with $\card T=n$ that is in convex position. There is no emptiness condition: points
of $S\setminus T$ may lie in $\conv T$. For $K\in\N$ we say that $S$ is \emph{$K$-free} if
$\neg\,\mathrm{HasConvexNGon}\,K\,S$.
\end{definition}

\begin{remark}\label{fd:rem:ngon}
$\mathrm{HasConvexNGon}\,0\,S$ holds for every $S$ (Lemma~\ref{fd:lem:small}(a)), so no set is
$0$-free, and a hypothesis ``$S$ is $0$-free'' is contradictory. If $S$ is $K$-free and $K\le K'$, then
$S$ is $K'$-free (contrapositive of Lemma~\ref{fd:lem:mono}(d)). If $\mathrm{HasConvexNGon}\,n\,S$, then
$n\le\card S$ (Lemma~\ref{fd:lem:ngon-basic}(a)).
\end{remark}

\begin{definition}[Forcing a convex $n$-gon; \leanref{ErdosSzekeres.ForcesConvexNGon}{ErdosSzekeres/Definitions.lean:25}]\label{fd:def:forces}\sloppy
For $N,n\in\N$, $\mathrm{ForcesConvexNGon}\,N\,n$ holds if every finite set $S$ of points that is in
general position and satisfies $N\le\card S$ satisfies $\mathrm{HasConvexNGon}\,n\,S$.
\end{definition}

\begin{definition}[Invariant sets; \leanref{ErdosSzekeres.LayerLemma.Invariant}{LAYERLEAN/LayerLemma.lean:105}, \leanref{ErdosSzekeres.SymReduction.Invariant}{SYMLEAN/SymmetricReduction.lean:271}, \leanref{ErdosSzekeres.OrbitConsistency.Invariant}{LEANEXT/orbit/OrbitConsistency.lean:47}]\label{fd:def:inv}\sloppy
For a map $f\colon\R^2\to\R^2$ and a finite set $S$ of points, $\mathrm{Invariant}\,f\,S$ holds if
$f(p)\in S$ for every $p\in S$; that is, $f$ maps $S$ \emph{into} $S$. For injective $f$ this is
equivalent to $f(S)=S$, and then also $f^{-1}(S)=\{p\in\R^2:f(p)\in S\}=S$
(Lemma~\ref{fd:lem:inv-image}). The three cited Lean definitions are textually identical.
\end{definition}

\begin{lemma}[Invariance and images; (a), (b) written step; (c) \leanref{ErdosSzekeres.LayerLemma.iterate_mem}{LAYERLEAN/LayerLemma.lean:52}]\label{fd:lem:inv-image}\sloppy
Let $X$ be a set, let $f\colon X\to X$ be a map and let $S$ be a finite subset of $X$.
\begin{enumerate}[label=(\alph*)]
\item If $X=\R^2$ and $f$ is injective, then $\mathrm{Invariant}\,f\,S$ holds if and only if $f(S)=S$.
\item If $X=\R^2$, $f$ is injective and $\mathrm{Invariant}\,f\,S$ holds, then $\{p\in\R^2:f(p)\in S\}=S$.
\item If $f(y)\in S$ for every $y\in S$, then $f^{[i]}(x)\in S$ for every $x\in S$ and every $i\in\N$.
\end{enumerate}
Part (c) is stated for an arbitrary set $X$, as in Lean, where $X$ is an arbitrary type; parts (a) and
(b) take $X=\R^2$ because $\mathrm{Invariant}$ is defined for maps of the plane.
\end{lemma}

\begin{proof}
(a) If $f(S)=S$, then $f(p)\in f(S)=S$ for every $p\in S$. Conversely, if $\mathrm{Invariant}\,f\,S$
holds, then $f(S)\subseteq S$; since $f$ is injective, $\card{f(S)}=\card S$, and a subset of a finite
set with the same number of elements is the whole set, so $f(S)=S$. (Lean uses this argument in the step
at line~165 of the proof of \lean{ErdosSzekeres.LayerLemma.map_mem_hullVerts},
\dirn{LAYERLEAN/LayerLemma.lean:162}.)

(b) If $p\in S$ then $f(p)\in S$ by invariance. If $f(p)\in S$, then by (a) $f(p)\in f(S)$, so
$f(p)=f(q)$ for some $q\in S$, and $p=q\in S$ by injectivity.

(c) Induction on $i$. For $i=0$, $f^{[0]}(x)=x\in S$. If $f^{[i]}(x)\in S$, then
$f^{[i+1]}(x)=f\bigl(f^{[i]}(x)\bigr)\in S$ by the hypothesis applied to $y=f^{[i]}(x)$. Here we used
$f^{[i+1]}=f\circ f^{[i]}$, which follows from the definition $f^{[i+1]}=f^{[i]}\circ f$ by induction on
$i$ ($f^{[1]}=f^{[0]}\circ f=f=f\circ f^{[0]}$, and $f^{[i+2]}=f^{[i+1]}\circ f=f\circ f^{[i]}\circ
f=f\circ f^{[i+1]}$; Mathlib's \lean{Function.iterate_succ_apply'}). This is the Lean proof.
\end{proof}

\begin{definition}[Rotation about a point; \leanref{ErdosSzekeres.LayerLemma.rotAbout}{LAYERLEAN/LayerLemma.lean:298}, \leanref{ErdosSzekeres.SymReduction.rotAbout}{SYMLEAN/SymmetricReduction.lean:400}, \leanref{ErdosSzekeres.OrbitConsistency.rotAbout}{LEANEXT/orbit/OrbitConsistency.lean:93}]\label{fd:def:rot}\sloppy
For $c\in\R^2$ and $\theta\in\R$, the map $\rho_{c,\theta}=\mathrm{rotAbout}\,c\,\theta\colon\R^2\to\R^2$
is
\[
  \rho_{c,\theta}(p)=\bigl(c_0+\cos\theta\,(p_0-c_0)-\sin\theta\,(p_1-c_1),\;
  c_1+\sin\theta\,(p_0-c_0)+\cos\theta\,(p_1-c_1)\bigr).
\]
In complex notation, identifying $p$ with $p_0+ip_1$ and $c$ with $c_0+ic_1$, this is
$z\mapsto c+e^{i\theta}(z-c)$ (Lemma~\ref{fd:lem:rot}(a)): the counterclockwise rotation about $c$ by
the angle $\theta$ (Lemma~\ref{fd:lem:rot}(e)). The three cited Lean definitions are textually
identical; a text search of the research records finds 19 definitions of \lean{rotAbout}, all with this
same body.
\end{definition}

\begin{remark}[Angles $2\pi/m$]\label{fd:rem:angles}
In later sections the angle is usually written \texttt{2 * Real.pi / m} in Lean with \lean{m} a natural
number, that is, $\theta=2\pi/(m:\R)$. For $m=0$ Lean's division gives $\theta=0$, and $\rho_{c,0}$ is
the identity (Lemma~\ref{fd:lem:rot}(c)). The propositions \lean{Sixfold66} and \lean{Sevenfold70}
(Definition~\ref{ll:def:66-70}) write their angles as \texttt{2 * Real.pi / 6} and
\texttt{2 * Real.pi / 7} with real numerals; since the real numbers $6$ and $7$ are the casts of the
natural numbers $6$ and $7$, these are the angles $2\pi/(m:\R)$ for $m=6$ and $m=7$.
\end{remark}

\begin{lemma}[Basic properties of rotations; (a) \leanref{ErdosSzekeres.LayerLemma.rotAbout_eq_crot}{LAYERLEAN/LayerLemma.lean:302}; (b) \leanref{ErdosSzekeres.LayerLemma.rotAff}{LAYERLEAN/LayerLemma.lean:311}; (c) \leanref{ErdosSzekeres.LayerLemma.rotAbout_iterate}{LAYERLEAN/LayerLemma.lean:465}; (d) \leanref{ErdosSzekeres.LayerLemma.rotAbout_injective}{LAYERLEAN/LayerLemma.lean:319}; (e), (f) written step]\label{fd:lem:rot}\sloppy
Let $c\in\R^2$ and $\theta,\alpha,\beta\in\R$. Write $\rho_\theta=\rho_{c,\theta}$ and
$R_\theta=\begin{pmatrix}\cos\theta&-\sin\theta\\ \sin\theta&\cos\theta\end{pmatrix}$.
\begin{enumerate}[label=(\alph*)]
\item For every $p\in\R^2$, the complex number $\rho_\theta(p)_0+i\,\rho_\theta(p)_1$ equals
  $\gamma+e^{i\theta}(z-\gamma)$, where $z=p_0+ip_1$ and $\gamma=c_0+ic_1$. (In Lean:
  $\mathrm{rotAbout}\,c\,\theta=\mathrm{crot}\,(\mathrm{toC}\,c)\,(\exp(\theta i))$, where
  $\mathrm{toC}\,p=p_0+ip_1$ and $\mathrm{crot}\,\gamma\,w\,p$ is the point with complex coordinate
  $\gamma+w(\mathrm{toC}\,p-\gamma)$.)
\item $\rho_\theta(p)=c+R_\theta(p-c)=R_\theta\,p+(c-R_\theta\,c)$ for every $p$; thus $\rho_\theta$ is an
  affine map with linear part $R_\theta$, that is, $\rho_\theta(p+v)=\rho_\theta(p)+R_\theta v$ for all
  $p,v\in\R^2$. In particular $\rho_\theta(c)=c$. (The Lean definition \lean{LayerLemma.rotAff} cited for (b)
  packages $\rho_\theta$ as an affine map with linear part $R_\theta$, its defining property being the
  identity $\rho_\theta(p+v)=\rho_\theta(p)+R_\theta v$. The matrix form
  $\rho_\theta(p)=c+R_\theta(p-c)=R_\theta\,p+(c-R_\theta\,c)$ of Definition~\ref{fd:def:rot} and the
  identity $\rho_\theta(c)=c$ are written steps here; Lean proves the latter in the step at line~330 of the
  proof of \lean{ErdosSzekeres.LayerLemma.invariant_erase_centre}, \dirn{LAYERLEAN/LayerLemma.lean:325}.)
\item $\rho_\alpha\circ\rho_\beta=\rho_{\alpha+\beta}$ and $\rho_0=\mathrm{id}$. Hence
  $\rho_\alpha^{[j]}=\rho_{(j:\R)\alpha}$ for every $j\in\N$. (The Lean declaration cited for (c) is this
  iterate formula; the composition law and $\rho_0=\mathrm{id}$ are written steps here.)
\item $\rho_\theta$ is a bijection of $\R^2$ with inverse $\rho_{-\theta}$; in particular it is
  injective. (The Lean declaration cited for (d) states injectivity; bijectivity and the inverse are a
  written step here.)
\item $\orient\bigl(c,p,\rho_\theta(p)\bigr)=\sin\theta\,\bigl((p_0-c_0)^2+(p_1-c_1)^2\bigr)$ for every
  $p$. Hence, if $0<\theta<\pi$ and $p\ne c$, then $\orient(c,p,\rho_\theta(p))>0$: the rotation turns
  counterclockwise.
\item $(\rho_\theta(p)_0-c_0)^2+(\rho_\theta(p)_1-c_1)^2=(p_0-c_0)^2+(p_1-c_1)^2$ for every $p$.
\end{enumerate}
\end{lemma}

\begin{proof}
We give direct proofs in coordinates; the Lean proofs of (c) and (d) pass through the complex form (a).
Put $u=p-c$. By definition,
\[
  \rho_\theta(p)-c=\bigl(\cos\theta\,u_0-\sin\theta\,u_1,\;\sin\theta\,u_0+\cos\theta\,u_1\bigr)
  =R_\theta\,u .
\]
(b) This is the displayed identity; $R_\theta(p-c)=R_\theta p-R_\theta c$ by linearity. Hence
$\rho_\theta(p+v)=c+R_\theta(p-c)+R_\theta v=\rho_\theta(p)+R_\theta v$, and $p=c$ gives $u=0$ and
$\rho_\theta(c)=c$.

(a) By Euler's formula $e^{i\theta}=\cos\theta+i\sin\theta$, and
$(\cos\theta+i\sin\theta)(u_0+iu_1)=(\cos\theta\,u_0-\sin\theta\,u_1)+i(\sin\theta\,u_0+\cos\theta\,u_1)$.
Adding $\gamma=c_0+ic_1$ gives the two coordinates of $\rho_\theta(p)$.

(c) By (b), $\rho_\alpha(\rho_\beta(p))=c+R_\alpha R_\beta(p-c)$. By the addition formulas
$\cos(\alpha+\beta)=\cos\alpha\cos\beta-\sin\alpha\sin\beta$ and
$\sin(\alpha+\beta)=\sin\alpha\cos\beta+\cos\alpha\sin\beta$,
\[
  R_\alpha R_\beta=
  \begin{pmatrix}\cos\alpha\cos\beta-\sin\alpha\sin\beta&-(\sin\alpha\cos\beta+\cos\alpha\sin\beta)\\
  \sin\alpha\cos\beta+\cos\alpha\sin\beta&\cos\alpha\cos\beta-\sin\alpha\sin\beta\end{pmatrix}
  =R_{\alpha+\beta},
\]
so $\rho_\alpha\circ\rho_\beta=\rho_{\alpha+\beta}$. Since $R_0$ is the identity matrix, $\rho_0=\mathrm{id}$.
Now induct on $j$: $\rho_\alpha^{[0]}=\mathrm{id}=\rho_0=\rho_{(0:\R)\alpha}$, and
$\rho_\alpha^{[j+1]}=\rho_\alpha^{[j]}\circ\rho_\alpha=\rho_{(j:\R)\alpha}\circ\rho_\alpha
=\rho_{(j:\R)\alpha+\alpha}=\rho_{(j+1:\R)\alpha}$.

(d) By (c), $\rho_{-\theta}\circ\rho_\theta=\rho_0=\mathrm{id}$ and $\rho_\theta\circ\rho_{-\theta}=\mathrm{id}$.

(e) With $u=p-c$ as above,
\[
  \orient(c,p,\rho_\theta(p))=u_0(\sin\theta\,u_0+\cos\theta\,u_1)-u_1(\cos\theta\,u_0-\sin\theta\,u_1)
  =\sin\theta\,(u_0^2+u_1^2).
\]
If $0<\theta<\pi$ then $\sin\theta>0$, and if $p\ne c$ then $u_0^2+u_1^2>0$.

(f) $(\cos\theta\,u_0-\sin\theta\,u_1)^2+(\sin\theta\,u_0+\cos\theta\,u_1)^2
=(\cos^2\theta+\sin^2\theta)(u_0^2+u_1^2)=u_0^2+u_1^2$, the cross terms cancelling.
\end{proof}

%======================================================================
\subsection{Facts}\label{fd:sec-facts}

\begin{lemma}[Symmetries of the orientation; (a) \leanref{ErdosSzekeres.orient_swap}{ErdosSzekeres/Orientation.lean:12}; (b) \leanref{ErdosSzekeres.orient_cycle}{ErdosSzekeres/Orientation.lean:15}; (c) \leanref{ErdosSzekeres.orient_translate}{ErdosSzekeres/Orientation.lean:18}; (d), (e) written step]\label{fd:lem:orient}\sloppy
Let $R$ be a commutative ring and let $a,b,c,v\in R^2$.
\begin{enumerate}[label=(\alph*)]
\item $\orient(a,c,b)=-\orient(a,b,c)$.
\item $\orient(b,c,a)=\orient(a,b,c)$. Consequently $\orient(a,b,c)=\orient(b,c,a)=\orient(c,a,b)$.
\item $\orient(a+v,b+v,c+v)=\orient(a,b,c)$.
\item For all $p_1,p_2,p_3\in R^2$ and every permutation $\pi$ of $\{1,2,3\}$,
  $\orient(p_{\pi(1)},p_{\pi(2)},p_{\pi(3)})=\sgn(\pi)\,\orient(p_1,p_2,p_3)$.
\item $\orient(a,b,c)=0$ whenever two of $a,b,c$ are equal. Consequently, if $\orient(a,b,c)\ne0$, then
  $a,b,c$ are pairwise distinct.
\end{enumerate}
\end{lemma}

\begin{proof}
The Lean proofs of (a)--(c) expand both sides and compare polynomials (\texttt{unfold orient; ring} for
(a) and (b), \texttt{simp only [orient, Pi.add\_apply]; ring} for (c)); we organise the same
computation. For $x,y\in R^2$ put $[x,y]=x_0y_1-x_1y_0$. Then $[\cdot,\cdot]$ is
bilinear, $[y,x]=-[x,y]$ and $[x,x]=0$. Expanding the definition (the terms $a_0a_1$ cancel),
\[
  \orient(a,b,c)=a_0b_1-a_1b_0+b_0c_1-b_1c_0+c_0a_1-c_1a_0=[a,b]+[b,c]+[c,a].
\]
(b) $\orient(b,c,a)=[b,c]+[c,a]+[a,b]=\orient(a,b,c)$. Applying this twice gives the second statement.

(a) $\orient(a,c,b)=[a,c]+[c,b]+[b,a]=-[c,a]-[b,c]-[a,b]=-\orient(a,b,c)$.

(c) The defining formula involves only the differences $b-a$ and $c-a$, and
$(b+v)-(a+v)=b-a$, $(c+v)-(a+v)=c-a$.

(d) The permutation group of $\{1,2,3\}$ has six elements. The identity and the two $3$-cycles have sign
$+1$; they give $\orient(p_1,p_2,p_3)$, $\orient(p_2,p_3,p_1)$ and $\orient(p_3,p_1,p_2)$, which are equal
by (b). The three transpositions have sign $-1$; they give $\orient(p_1,p_3,p_2)$, $\orient(p_2,p_1,p_3)$
and $\orient(p_3,p_2,p_1)$. The first equals $-\orient(p_1,p_2,p_3)$ by (a). By (b) applied to the triple
$(p_2,p_1,p_3)$ we get $\orient(p_1,p_3,p_2)=\orient(p_2,p_1,p_3)$, and by (b) applied to
$(p_3,p_2,p_1)$ we get $\orient(p_2,p_1,p_3)=\orient(p_3,p_2,p_1)$; so all three equal
$-\orient(p_1,p_2,p_3)$.

(e) If $a=b$ the first factors $b_0-a_0$ and $b_1-a_1$ vanish; if $a=c$ the factors $c_1-a_1$ and
$c_0-a_0$ vanish; if $b=c$ the two products $(b_0-a_0)(b_1-a_1)$ and $(b_1-a_1)(b_0-a_0)$ are equal. In
each case $\orient(a,b,c)=0$.
\end{proof}

\begin{remark}[The real case of (d) in Lean]\label{fd:rem:orient-lean}\sloppy
For $R=\R$, part (d) of Lemma~\ref{fd:lem:orient} is the Lean theorem
\lean{ErdosSzekeres.OrbitConsistency.orient3_perm} (\dirn{LEANEXT/orbit/OrbitConsistency.lean:178};
Lemma~\ref{oc:lem:perm}). It is stated for ordered triples $t\colon\mathrm{Fin}\,3\to\R^2$, with
$\orient_3(t)=\orient(t(0),t(1),t(2))$ (Lean \lean{orient3},
\dirn{LEANEXT/orbit/OrbitConsistency.lean:168}), and says that
$\orient_3(t\circ\sigma)=\sgn(\sigma)\,\orient_3(t)$ for every permutation $\sigma$ of
$\mathrm{Fin}\,3=\{0,1,2\}$ and every $t$. Put $t(i)=p_{i+1}$ and $\sigma(i)=\pi(i+1)-1$; then
$\sigma=s^{-1}\circ\pi\circ s$ for the bijection $s(i)=i+1$ from $\{0,1,2\}$ to $\{1,2,3\}$, so $\sigma$
has the same cycle type and hence the same sign as $\pi$, and $(t\circ\sigma)(i)=p_{\pi(i+1)}$. So
$\orient_3(t\circ\sigma)=\orient(p_{\pi(1)},p_{\pi(2)},p_{\pi(3)})$ and $\orient_3(t)=\orient(p_1,p_2,p_3)$,
and the Lean statement is (d) for $R=\R$. Parts (d) for a general commutative ring and (e) are written
steps.
\end{remark}

\begin{lemma}[Orientation under affine maps; \leanref{ErdosSzekeres.SymReduction.orient_affine}{SYMLEAN/AffineReduction.lean:1241}; (i) \leanref{ErdosSzekeres.OrbitConsistency.orient_rotAbout_inv}{LEANEXT/orbit/OrbitConsistency.lean:117}; (ii) written step]\label{fd:lem:affine-orient}\sloppy
Let $g\colon\R^2\to\R^2$ be a map, let $A$ be a real $2\times2$ matrix and $v\in\R^2$, and suppose that
$g(x)=Ax+v$ for every $x\in\R^2$. Then for all $a,b,c\in\R^2$
\[
  \orient\bigl(g(a),g(b),g(c)\bigr)=\det(A)\,\orient(a,b,c).
\]
(In Lean, $g$ is an arbitrary map with the hypothesis that $g(p)=L\,p+w$ for every $p$, where $L$ is a
real $2\times2$ matrix and $w\in\R^2$ (Lean's \lean{b}), and the conclusion is
$\orient(g(x),g(y),g(z))=\det(L)\,\orient(x,y,z)$ for all $x,y,z$; this is the statement above with
$L=A$ and $w=v$.)
In particular:
\begin{enumerate}[label=(\roman*)]
\item (rotations have determinant $1$) for all $o\in\R^2$, $\theta\in\R$ and $a,b,d\in\R^2$,
  $\orient(\rho_{o,\theta}(a),\rho_{o,\theta}(b),\rho_{o,\theta}(d))=\orient(a,b,d)$, because the linear part
  $R_\theta$ of $\rho_{o,\theta}$ (Lemma~\ref{fd:lem:rot}(b)) has $\det R_\theta=\cos^2\theta+\sin^2\theta=1$;
\item (reflections have determinant $-1$) for $o\in\R^2$ and $\varphi\in\R$ let
  $F_\varphi=\begin{pmatrix}\cos2\varphi&\sin2\varphi\\ \sin2\varphi&-\cos2\varphi\end{pmatrix}$ and
  $\sigma(x)=o+F_\varphi(x-o)$, the reflection in the line through $o$ with direction
  $u=(\cos\varphi,\sin\varphi)$ (indeed $F_\varphi u=u$ and $F_\varphi u^\perp=-u^\perp$ for
  $u^\perp=(-\sin\varphi,\cos\varphi)$). Then
  $\orient(\sigma(a),\sigma(b),\sigma(d))=-\orient(a,b,d)$ for all $a,b,d$, because
  $\det F_\varphi=-\cos^22\varphi-\sin^22\varphi=-1$. For $o=0$ and $\varphi=0$, $\sigma$ is the reflection
  $(p_0,p_1)\mapsto(p_0,-p_1)$ in the first coordinate axis.
\end{enumerate}
\end{lemma}

\begin{proof}
By Definition~\ref{fd:def:orient}, $\orient(a,b,c)=\det(b-a\mid c-a)$. Since $g(b)-g(a)=A(b-a)$ and
$g(c)-g(a)=A(c-a)$, the matrix with columns $g(b)-g(a)$ and $g(c)-g(a)$ is the product $A\,(b-a\mid c-a)$,
and the determinant of a product of $2\times2$ matrices is the product of the determinants (a polynomial
identity in the eight entries, checked by expanding both sides). This proves the formula. (The Lean proof
substitutes $g(p)=Lp+w$, writes $\det L=L_{00}L_{11}-L_{01}L_{10}$, expands both sides in coordinates
and compares polynomials, which is the same computation.)

(i) $\rho_{o,\theta}(x)=R_\theta x+(o-R_\theta o)$ by Lemma~\ref{fd:lem:rot}(b), and
$\det R_\theta=\cos\theta\cdot\cos\theta-(-\sin\theta)\sin\theta=1$.

(ii) $\sigma(x)=F_\varphi x+(o-F_\varphi o)$ has linear part $F_\varphi$, and
$\det F_\varphi=-\cos^22\varphi-\sin^22\varphi=-1$. Using $\cos2\varphi\cos\varphi+\sin2\varphi\sin\varphi
=\cos\varphi$ and $\sin2\varphi\cos\varphi-\cos2\varphi\sin\varphi=\sin\varphi$ (the subtraction formulas for
$2\varphi-\varphi$) one gets $F_\varphi u=u$ and $F_\varphi u^\perp=(\sin\varphi,-\cos\varphi)=-u^\perp$, so
$\sigma$ fixes the line $o+\R u$ pointwise and reverses the perpendicular direction. For $\varphi=0$,
$F_0=\mathrm{diag}(1,-1)$.
\end{proof}

\begin{remark}[The formula and its special cases in Lean]\label{fd:rem:affine-orient-lean}\sloppy
The general formula is the Lean theorem \lean{ErdosSzekeres.SymReduction.orient_affine}, cited in the
head; a verbatim copy, with the same name, is at \dirn{SYMLEAN/mirror/ReductionThreeClasses.lean:1267}.
Part (i) is a separate Lean theorem, \lean{ErdosSzekeres.OrbitConsistency.orient_rotAbout_inv}, also
cited in the head (Lemma~\ref{oc:lem:rot-inv}). It is stated for the copy
\lean{ErdosSzekeres.OrbitConsistency.rotAbout} of the rotation of Definition~\ref{fd:def:rot}, and it
calls the centre $c$ instead of $o$. Part (ii) is a written step. Its case $o=0$, $\varphi=0$, the sign
change under $(p_0,p_1)\mapsto(p_0,-p_1)$, is the Lean theorem \lean{ErdosSzekeres.Suk.orient_reflectY}
(\dirn{ErdosSzekeres/Suk/Reflection.lean:47}).
\end{remark}

\begin{lemma}[Collinearity through a point; Mathlib \lean{collinear_iff_of_mem}; written step]\label{fd:lem:collinear-mem}
Let $X\subseteq\R^2$ and $p^*\in X$. Then $X$ is collinear if and only if there is $v\in\R^2$ such that
every $p\in X$ can be written $p=p^*+r\,v$ with $r\in\R$.
\end{lemma}

\begin{proof}
Let $W$ be the linear span of $\{p-q:p,q\in X\}$. Since $p^*\in X$, $W$ is also the span of
$\{p-p^*:p\in X\}$: these vectors are among the differences, and conversely
$p-q=(p-p^*)-(q-p^*)$. If $X$ is collinear, then $\dim W\le1$, so $W\subseteq\R v$ for some $v$ (take
$v=0$ if $W=0$, and a basis vector of $W$ if $\dim W=1$); then $p-p^*\in W$ gives $p=p^*+rv$ for some
$r$. Conversely, if every $p-p^*$ lies in $\R v$, then $W\subseteq\R v$ and $\dim W\le\dim\R v\le1$.
\end{proof}

\begin{lemma}[Orientation and collinearity; (a) \leanref{ErdosSzekeres.orient_eq_zero_iff_collinear}{ErdosSzekeres/OrientationGeometry.lean:36}; (b) \leanref{ErdosSzekeres.generalPosition_iff_orient_ne_zero}{ErdosSzekeres/OrientationGeometry.lean:66}]\label{fd:lem:collinear}\sloppy\tolerance=10000
The orientation detects collinearity, and hence general position.
\begin{enumerate}[label=(\alph*)]
\item For all $a,b,c\in\R^2$: $\orient(a,b,c)=0$ if and only if the set $\{a,b,c\}$ is collinear. This
  includes degenerate triples, in which two or three of the points coincide.
\item A finite set $S$ of points is in general position if and only if $\orient(a,b,c)\ne0$ for all
  $a,b,c\in S$ with $a\ne b$, $a\ne c$ and $b\ne c$.
\end{enumerate}
\end{lemma}

\begin{proof}
We follow the Lean proof. (a) We use Lemma~\ref{fd:lem:collinear-mem} with $X=\{a,b,c\}$ and $p^*=a$.

Suppose $\{a,b,c\}$ is collinear. Then there are $v$ and reals $r,s$ with $b=a+rv$ and $c=a+sv$, so
$\orient(a,b,c)=(rv_0)(sv_1)-(rv_1)(sv_0)=0$.

Suppose $\orient(a,b,c)=0$. If $b=a$, put $v=c-a$; then $a=a+0\cdot v$, $b=a+0\cdot v$ and $c=a+1\cdot v$, so
$\{a,b,c\}$ is collinear. If $b\ne a$, put $u=b-a\ne0$ and $w=c-a$; then
$u_0w_1-u_1w_0=\orient(a,b,c)=0$. We find $r$ with $w=ru$. If $u_0\ne0$, let $r=w_0/u_0$; then $ru_0=w_0$
and $ru_1=w_0u_1/u_0=u_0w_1/u_0=w_1$. If $u_0=0$, then $u_1\ne0$ and $0=u_0w_1-u_1w_0=-u_1w_0$ gives
$w_0=0$; let $r=w_1/u_1$, so that $ru_0=0=w_0$ and $ru_1=w_1$. Now $a=a+0\cdot u$, $b=a+1\cdot u$ and
$c=a+r\,u$, so $\{a,b,c\}$ is collinear.

(b) By Definition~\ref{fd:def:gp}, $S$ is in general position iff $\{a,b,c\}$ is not collinear for all
pairwise distinct $a,b,c\in S$; by (a) this says $\orient(a,b,c)\ne0$ for all such triples.
\end{proof}

\begin{lemma}[Convex hulls; written step]\label{fd:lem:hull}
Let $Y,Y'\subseteq\R^2$.
\begin{enumerate}[label=(\alph*)]
\item $Y\subseteq\conv Y$; $\conv Y$ is convex; and $\conv Y\subseteq K$ for every convex set $K$ with
  $Y\subseteq K$. If $Y\subseteq Y'$ then $\conv Y\subseteq\conv Y'$. Moreover $\conv\emptyset=\emptyset$
  and $\conv\{y\}=\{y\}$.
\item $\conv Y$ is the set of all convex combinations of finitely many points of $Y$.
\item If $F=\{y_1,\dots,y_m\}$ is a finite set of $m\ge1$ distinct points, then
  $\conv F=\bigl\{\sum_{i=1}^m\lambda_iy_i:\lambda_i\ge0,\ \sum_i\lambda_i=1\bigr\}$. In particular
  $\conv\{y,z\}=\{\lambda y+(1-\lambda)z:0\le\lambda\le1\}$.
\item For $w\in\R^2$ and $\alpha\in\R$, the open half-plane $\{z\in\R^2:w_0z_0+w_1z_1<\alpha\}$ is
  convex.
\end{enumerate}
\end{lemma}

\begin{proof}
(a) An intersection of convex sets is convex, so $\conv Y$ is convex; it contains $Y$; and it is contained
in each convex $K\supseteq Y$, being the intersection of all of them. If $Y\subseteq Y'$, every convex
set containing $Y'$ contains $Y$, so $\conv Y\subseteq\conv Y'$. The empty set and $\{y\}$ are convex
(for $\{y\}$: $\lambda y+\mu y=y$ when $\lambda+\mu=1$), which gives the last two claims.

(b) Let $C$ be the set of convex combinations of finitely many points of $Y$. Every $y\in Y$ is the convex
combination $1\cdot y$, so $Y\subseteq C$. $C$ is convex: if $x=\sum_i\lambda_iy_i$ and
$x'=\sum_j\lambda'_jy'_j$ are in $C$ and $\lambda,\mu\ge0$ with $\lambda+\mu=1$, then
$\lambda x+\mu x'=\sum_i(\lambda\lambda_i)y_i+\sum_j(\mu\lambda'_j)y'_j$ has nonnegative coefficients
with sum $\lambda+\mu=1$. Hence $\conv Y\subseteq C$ by (a). Conversely, let $K\supseteq Y$ be convex; we
show by induction on $k$ that every convex combination $x=\sum_{i=1}^k\lambda_iy_i$ of points of $Y$
lies in $K$. For $k=1$, $x=y_1\in K$. For $k+1$ terms: if $\lambda_{k+1}=1$, all other coefficients
vanish and $x=y_{k+1}\in K$; otherwise $s=\sum_{i\le k}\lambda_i=1-\lambda_{k+1}>0$, the point
$x'=\sum_{i\le k}(\lambda_i/s)y_i$ is a convex combination of $k$ points of $Y$, so $x'\in K$, and
$x=s\,x'+\lambda_{k+1}y_{k+1}$ with $s+\lambda_{k+1}=1$ lies in $K$ by convexity. Taking $K=\conv Y$
gives $C\subseteq\conv Y$.

(c) By (b), $\conv F$ is the set of convex combinations of points of $F$. In such a combination we may
merge the terms that use the same point of $F$ (adding their coefficients) and add the missing points of
$F$ with coefficient $0$; this gives the stated form, and conversely every point of the stated form is a
convex combination. For $F=\{y,z\}$ with $y\ne z$ this is the segment; for $y=z$ both sides are $\{y\}$.

(d) If $w_0z_0+w_1z_1<\alpha$ and $w_0z'_0+w_1z'_1<\alpha$, and $\lambda,\mu\ge0$ with $\lambda+\mu=1$,
then $w_0(\lambda z_0+\mu z'_0)+w_1(\lambda z_1+\mu z'_1)<\lambda\alpha+\mu\alpha=\alpha$, because at least
one of $\lambda,\mu$ is positive and the other term is at most the corresponding multiple of $\alpha$.
\end{proof}

\begin{lemma}[Convex position via hulls; \leanref{ErdosSzekeres.convexPosition_iff}{ErdosSzekeres/Basic.lean:15}]\label{fd:lem:cp-hull}\sloppy
A finite set $S$ of points is in convex position if and only if $p\notin\conv(S\setminus\{p\})$ for
every $p\in S$.
\end{lemma}

\begin{proof}
In Lean this is Mathlib's \lean{convexIndependent_set_iff_notMem_convexHull_diff} applied to $S$; we give
a direct proof from the form of Definition~\ref{fd:def:cp}.

Suppose $S$ is in convex position and $p\in S$. If $p\in\conv(S\setminus\{p\})$, then the definition with
$T=S\setminus\{p\}$ and $x=p$ gives $p\in S\setminus\{p\}$, which is absurd.

Conversely, suppose $p\notin\conv(S\setminus\{p\})$ for every $p\in S$. Let $T\subseteq S$ and $x\in S$
with $x\in\conv T$, and suppose $x\notin T$. Then $T\subseteq S\setminus\{x\}$, so
$x\in\conv T\subseteq\conv(S\setminus\{x\})$ by Lemma~\ref{fd:lem:hull}(a), a contradiction. Hence
$x\in T$, and $S$ is in convex position.
\end{proof}

\begin{lemma}[Monotonicity; (a) \leanref{ErdosSzekeres.ConvexPosition.mono}{ErdosSzekeres/Basic.lean:11}; (b) \leanref{ErdosSzekeres.GeneralPosition.mono}{ErdosSzekeres/Basic.lean:6}; (c) \leanref{ErdosSzekeres.HasConvexNGon.mono}{ErdosSzekeres/Basic.lean:20}; (d) \leanref{MirrorOdd.hasConvexNGon_of_le}{MIRRORETV/lean/MirrorMain.lean:2182}]\label{fd:lem:mono}\sloppy
Let $s,t,S$ be finite sets of points and $n,n'\in\N$.
\begin{enumerate}[label=(\alph*)]
\item If $s$ is in convex position and $t\subseteq s$, then $t$ is in convex position.
\item If $s$ is in general position and $t\subseteq s$, then $t$ is in general position.
\item If $\mathrm{HasConvexNGon}\,n\,s$ and $s\subseteq t$, then $\mathrm{HasConvexNGon}\,n\,t$.
\item If $\mathrm{HasConvexNGon}\,n\,S$ and $n'\le n$, then $\mathrm{HasConvexNGon}\,n'\,S$.
\end{enumerate}
\end{lemma}

\begin{proof}
(a) Let $p\in t$. Then $t\setminus\{p\}\subseteq s\setminus\{p\}$, so
$\conv(t\setminus\{p\})\subseteq\conv(s\setminus\{p\})$ (Lemma~\ref{fd:lem:hull}(a)), and the latter does
not contain $p$ by Lemma~\ref{fd:lem:cp-hull}, since $p\in s$. By Lemma~\ref{fd:lem:cp-hull} again, $t$ is
in convex position. (Lean uses Mathlib's \lean{ConvexIndependent.mono}.)

(b) Pairwise distinct $a,b,c\in t$ are pairwise distinct points of $s$, so $\{a,b,c\}$ is not collinear.
This is the Lean proof.

(c) A witness $T\subseteq s$ with $\card T=n$ in convex position satisfies $T\subseteq t$. This is the
Lean proof.

(d) Let $T\subseteq S$ with $\card T=n$ be in convex position. Since $n'\le n=\card T$, there is
$T'\subseteq T$ with $\card{T'}=n'$; it is in convex position by (a), and $T'\subseteq S$. This is the
Lean proof (it applies Lemma~\ref{fd:lem:ngon-basic}(b) to $T$ and then part (c) above). The same declaration, with the same name and
proof, also appears at \dirn{SYMLEAN/mirror/ReductionThreeClasses.lean:3741}.
\end{proof}

\begin{lemma}[Sizes of convex subsets; (a) \leanref{ErdosSzekeres.HasConvexNGon.card_le}{ErdosSzekeres/Basic.lean:25}; (b) \leanref{ErdosSzekeres.ConvexPosition.hasConvexNGon}{ErdosSzekeres/Basic.lean:31}]\label{fd:lem:ngon-basic}\sloppy
Let $s$ be a finite set of points and $n\in\N$.
\begin{enumerate}[label=(\alph*)]
\item If $\mathrm{HasConvexNGon}\,n\,s$, then $n\le\card s$.
\item If $s$ is in convex position and $n\le\card s$, then $\mathrm{HasConvexNGon}\,n\,s$.
\end{enumerate}
\end{lemma}

\begin{proof}
(a) A witness $T\subseteq s$ has $n=\card T\le\card s$. (b) Choose $T\subseteq s$ with $\card T=n$; it is
in convex position by Lemma~\ref{fd:lem:mono}(a). Both are the Lean proofs.
\end{proof}

\begin{lemma}[Forcing; (a) \leanref{ErdosSzekeres.ForcesConvexNGon.mono}{ErdosSzekeres/Basic.lean:47}; (b) \leanref{ErdosSzekeres.forcesConvexNGon_iff_exact}{ErdosSzekeres/Basic.lean:37}]\label{fd:lem:forces-basic}\sloppy
Let $N,M,n\in\N$.
\begin{enumerate}[label=(\alph*)]
\item If $\mathrm{ForcesConvexNGon}\,N\,n$ and $N\le M$, then $\mathrm{ForcesConvexNGon}\,M\,n$.
\item $\mathrm{ForcesConvexNGon}\,N\,n$ holds if and only if every finite set $S$ of points in general
  position with $\card S=N$ satisfies $\mathrm{HasConvexNGon}\,n\,S$.
\end{enumerate}
\end{lemma}

\begin{proof}
(a) If $S$ is in general position and $M\le\card S$, then $N\le\card S$. (b) The forward direction is the
case $\card S=N$. Conversely, let $S$ be in general position with $N\le\card S$. Choose $T\subseteq S$ with
$\card T=N$. By Lemma~\ref{fd:lem:mono}(b) $T$ is in general position, so $\mathrm{HasConvexNGon}\,n\,T$
by hypothesis, and $\mathrm{HasConvexNGon}\,n\,S$ by Lemma~\ref{fd:lem:mono}(c). Both are the Lean
proofs.
\end{proof}

\begin{lemma}[At most three points; written step, the last assertion \leanref{ErdosSzekeres.convexPosition_of_card_three}{ErdosSzekeres/Basic.lean:58}]\label{fd:lem:small}\sloppy
Sets of at most three points behave as follows.
\begin{enumerate}[label=(\alph*)]
\item For every finite set $S$ of points, $\mathrm{HasConvexNGon}\,0\,S$ holds.
\item Every finite set of at most two points is in convex position.
\item Let $a,b,c$ be pairwise distinct points. Then $\{a,b,c\}$ is in convex position if and only if
  $\{a,b,c\}$ is not collinear. In particular, if a finite set $s$ is in general position and
  $\card s=3$, then $s$ is in convex position.
\end{enumerate}
\end{lemma}

\begin{proof}
(a) $T=\emptyset$ is a subset of $S$ with $0$ elements, and it is in convex position by
Lemma~\ref{fd:lem:cp-hull} (there is nothing to check).

(b) We use Lemma~\ref{fd:lem:cp-hull} and Lemma~\ref{fd:lem:hull}(a). For $\emptyset$ there is nothing to
check; for $\{x\}$, $x\notin\conv\emptyset=\emptyset$; for $\{x,y\}$ with $x\ne y$,
$x\notin\conv\{y\}=\{y\}$ and $y\notin\conv\{x\}=\{x\}$.

(c) Since $a,b,c$ are pairwise distinct, $\{a,b,c\}\setminus\{a\}=\{b,c\}$, and similarly for $b$ and
$c$.

Suppose $\{a,b,c\}$ is not collinear. If $a\in\conv\{b,c\}$, then by Lemma~\ref{fd:lem:hull}(c)
$a=\lambda b+(1-\lambda)c$ for some $\lambda\in[0,1]$; with $v=b-c$ we get $a=c+\lambda v$, $b=c+1\cdot v$
and $c=c+0\cdot v$, so $\{a,b,c\}$ is collinear by Lemma~\ref{fd:lem:collinear-mem} (with $p^*=c$), a
contradiction. The same argument with the roles of $a,b,c$ exchanged shows $b\notin\conv\{a,c\}$ and
$c\notin\conv\{a,b\}$. By Lemma~\ref{fd:lem:cp-hull}, $\{a,b,c\}$ is in convex position.

Suppose $\{a,b,c\}$ is collinear. By Lemma~\ref{fd:lem:collinear-mem} with $p^*=a$ there are $v$ and
reals $r_b,r_c$ with $b=a+r_bv$ and $c=a+r_cv$; put $\tau(a)=0$, $\tau(b)=r_b$, $\tau(c)=r_c$, so that
every $x\in\{a,b,c\}$ equals $a+\tau(x)v$. Since the points are pairwise distinct, the three reals
$0,r_b,r_c$ are pairwise distinct (equal parameters would give equal points). Let $x$ be the point whose
parameter $\tau(x)$ is the middle one of the three, and let
$y,z$ be the other two points, with $\tau(y)<\tau(x)<\tau(z)$. With
$\lambda=(\tau(z)-\tau(x))/(\tau(z)-\tau(y))\in(0,1)$ we have
$\tau(x)=\lambda\tau(y)+(1-\lambda)\tau(z)$, hence $x=\lambda y+(1-\lambda)z\in\conv\{y,z\}
=\conv(\{a,b,c\}\setminus\{x\})$, and $\{a,b,c\}$ is not in convex position by
Lemma~\ref{fd:lem:cp-hull}.

Finally, if $s$ is in general position and $\card s=3$, then $s=\{a,b,c\}$ with $a,b,c$ pairwise distinct
and $\{a,b,c\}$ not collinear, so $s$ is in convex position. (The Lean proof of this last assertion
derives affine independence of $a,b,c$ from non-collinearity and uses that affinely independent families
are convex independent; we argued directly.)
\end{proof}

\begin{lemma}[Injective affine maps; (a) Mathlib \lean{AffineMap.image_convexHull}; (b$'$) \leanref{ErdosSzekeres.SymReduction.convexPosition_image_affine}{SYMLEAN/AffineReduction.lean:1263}; (c$'$) \leanref{ErdosSzekeres.SymReduction.generalPosition_image_affine}{SYMLEAN/AffineReduction.lean:1278}, \leanref{ErdosSzekeres.SymReduction.generalPosition_image_affineEquiv}{SYMLEAN/AffineReduction.lean:1290}; (d$'$) \leanref{ErdosSzekeres.SymReduction.hasConvexNGon_image_affine}{SYMLEAN/AffineReduction.lean:1298}; (e$'$) \leanref{ErdosSzekeres.SymReduction.injective_of_affine}{SYMLEAN/AffineReduction.lean:1305}; the rest written step]\label{fd:lem:affine-cp}\raggedright
Let $g\colon\R^2\to\R^2$ be an affine map, $g(x)=Ax+v$.
\begin{enumerate}[label=(\alph*)]
\item For every $Y\subseteq\R^2$, $g(\conv Y)=\conv\bigl(g(Y)\bigr)$.
\end{enumerate}
Suppose moreover that $g$ is injective. Then $\det A\ne0$, and $g$ is a bijection of $\R^2$ whose inverse
$y\mapsto A^{-1}(y-v)$ is affine. For every finite set $S$ of points and every $n\in\N$ (with $g(S)$ the
image set, $\card{g(S)}=\card S$):
\begin{enumerate}[label=(\alph*),resume]
\item $S$ is in convex position if and only if $g(S)$ is in convex position;
\item $S$ is in general position if and only if $g(S)$ is in general position;
\item $\mathrm{HasConvexNGon}\,n\,S$ holds if and only if $\mathrm{HasConvexNGon}\,n\,g(S)$ holds.
\end{enumerate}
The forward directions of (b)--(d), and the converse of the implication ``$g$ injective $\Rightarrow$
$\det A\ne0$'', are Lean theorems in the following forms. Here $S$ is a finite set of points, $n\in\N$,
$f$ is an affine map of $\R^2$ in the sense of Mathlib (\lean{AffineMap}; Definition~\ref{fd:def:conv}),
and $h\colon\R^2\to\R^2$ is a map with $h(p)=Lp+w$ for every $p\in\R^2$, where $L$ is a real $2\times2$
matrix and $w\in\R^2$ (Lean's \lean{b}).
\begin{enumerate}[label=(\alph*$'$),start=2]
\item If $S$ is in convex position and $f$ is injective, then $f(S)$ is in convex position.
\item If $\det L\ne0$ and $S$ is in general position, then $h(S)$ is in general position. If $S$ is in
  general position and $\phi$ is an affine equivalence of $\R^2$ in the sense of Mathlib
  (\lean{AffineEquiv}; Definition~\ref{fd:def:conv}), then $\phi(S)$ is in general position.
\item If $\mathrm{HasConvexNGon}\,n\,S$ and $f$ is injective, then $\mathrm{HasConvexNGon}\,n\,f(S)$.
\item If $\det L\ne0$, then $h$ is injective.
\end{enumerate}
\end{lemma}

\begin{proof}
For $x,x'\in\R^2$ and reals $\lambda,\mu$ with $\lambda+\mu=1$,
\begin{equation*}
  g(\lambda x+\mu x')=\lambda Ax+\mu Ax'+(\lambda+\mu)v=\lambda g(x)+\mu g(x'). \tag{$*$}
\end{equation*}

(a) By $(*)$, $g(\conv Y)$ is convex: any two of its points have the form $g(x),g(x')$ with
$x,x'\in\conv Y$, and for
$\lambda,\mu\ge0$ with $\lambda+\mu=1$, $\lambda g(x)+\mu g(x')=g(\lambda x+\mu x')$ with
$\lambda x+\mu x'\in\conv Y$. It contains $g(Y)$, so $\conv(g(Y))\subseteq g(\conv Y)$ by
Lemma~\ref{fd:lem:hull}(a). Conversely, the set $K=\{x:g(x)\in\conv(g(Y))\}$ is convex by $(*)$ and
contains $Y$, so $\conv Y\subseteq K$, that is, $g(\conv Y)\subseteq\conv(g(Y))$. (This is the proof of
Mathlib's \lean{AffineMap.image_convexHull}.)

Now let $g$ be injective, and write $A=\begin{pmatrix}\alpha&\beta\\ \gamma&\delta\end{pmatrix}$. If
$\det A=\alpha\delta-\beta\gamma=0$, there is $x\ne0$ with $Ax=0$: if $(\alpha,\beta)\ne(0,0)$ take
$x=(\beta,-\alpha)$, for which $Ax=(\alpha\beta-\beta\alpha,\ \gamma\beta-\delta\alpha)=(0,0)$; if
$\alpha=\beta=0$ take $x=(\delta,-\gamma)$ when $(\gamma,\delta)\ne(0,0)$, for which
$Ax=(0,\gamma\delta-\delta\gamma)=(0,0)$, and any $x\ne0$ when $A=0$. Then $g(x)=v=g(0)$ with $x\ne0$,
contradicting injectivity. So $\det A\ne0$, $A$ has the inverse
$A^{-1}=(\det A)^{-1}\begin{pmatrix}\delta&-\beta\\ -\gamma&\alpha\end{pmatrix}$, and $g$ is bijective with
inverse $y\mapsto A^{-1}(y-v)$, which is affine. Injectivity gives $\card{g(S)}=\card S$.

(b) The points of $g(S)$ are the $g(p)$ with $p\in S$, and $g(S)\setminus\{g(p)\}=g(S\setminus\{p\})$ by
injectivity. By (a), $\conv\bigl(g(S)\setminus\{g(p)\}\bigr)=g\bigl(\conv(S\setminus\{p\})\bigr)$, and by
injectivity $g(p)$ belongs to this set if and only if $p\in\conv(S\setminus\{p\})$. Lemma~\ref{fd:lem:cp-hull}
applied to $S$ and to $g(S)$ gives (b).

(c) Every triple of points of $g(S)$ has the form $(g(a),g(b),g(c))$ with $a,b,c\in S$; by injectivity
$g(a),g(b),g(c)$ are pairwise distinct if and only if $a,b,c$ are; and
$\orient(g(a),g(b),g(c))=\det(A)\,\orient(a,b,c)$ (Lemma~\ref{fd:lem:affine-orient}) with $\det A\ne0$, so
one orientation vanishes if and only if the other does. Lemma~\ref{fd:lem:collinear}(b) gives (c).

(d) If $T\subseteq S$, $\card T=n$, is in convex position, then $g(T)\subseteq g(S)$, $\card{g(T)}=n$, and
$g(T)$ is in convex position by (b). Conversely, let $T'\subseteq g(S)$ with $\card{T'}=n$ be in convex
position, and let $T=\{p\in S:g(p)\in T'\}$. Since every point of $T'$ is $g(p)$ for some $p\in S$, we have
$g(T)=T'$; hence $\card T=\card{T'}=n$ by injectivity, $T$ is in convex position by (b), and $T\subseteq S$.

(b$'$)--(e$'$). By Definition~\ref{fd:def:conv}, a Mathlib affine map $f$ of $\R^2$ has the form
$f(x)=A_fx+f(0)$, where $A_f$ is the matrix of its linear part (in Lean:
\lean{ErdosSzekeres.SymReduction.affineMap_eq_mulVec}, \dirn{SYMLEAN/AffineReduction.lean:1284}). So
(b$'$) and (d$'$) are the forward directions of (b) and (d) for $g=f$. (e$'$): if $h(p)=h(q)$, then
$L(p-q)=(Lp+w)-(Lq+w)=h(p)-h(q)=0$. Write $L=\begin{pmatrix}\alpha&\beta\\ \gamma&\delta\end{pmatrix}$; the
matrix $L^\sharp=\begin{pmatrix}\delta&-\beta\\ -\gamma&\alpha\end{pmatrix}$ satisfies
$L^\sharp L=(\det L)\,I$ (multiply out), so $(\det L)(p-q)=L^\sharp L(p-q)=0$, and $p=q$ because
$\det L\ne0$. (c$'$): if $\det L\ne0$, then $h$ is an affine map with linear part $L$, and it is injective
by (e$'$), so the forward direction of (c) applies to $g=h$. A Mathlib affine equivalence $\phi$ of
$\R^2$ is an injective affine map (Definition~\ref{fd:def:conv}), so the forward direction of (c) applies
to $g=\phi$ as well. (The Lean proofs follow the same routes:
\lean{SymReduction.convexPosition_image_affine} uses \lean{AffineMap.image_convexHull} and injectivity as
in (b); \lean{SymReduction.generalPosition_image_affine} uses \lean{SymReduction.orient_affine}, that is,
Lemma~\ref{fd:lem:affine-orient}; \lean{SymReduction.generalPosition_image_affineEquiv} derives
$\det\ne0$ for the matrix of the linear part of $\phi$ from the invertibility of that linear part.)
\end{proof}

\begin{remark}[Special cases in the main library]\label{fd:rem:affine-cp-lean}\sloppy
Besides the general forms (b$'$)--(e$'$), the main library proves special cases of
Lemma~\ref{fd:lem:affine-cp} in \dirn{ErdosSzekeres/Suk/Reflection.lean}.
\lean{ErdosSzekeres.HasConvexNGon.linearEquiv} (\dirn{ErdosSzekeres/Suk/Reflection.lean:6}) transports
$\mathrm{HasConvexNGon}\,n\,S$ to $\mathrm{HasConvexNGon}\,n\,e(S)$ for a linear bijection $e$ of the
plane (Mathlib \lean{LinearEquiv}): the forward direction of (d) for $g=e$. The reflection
$\sigma(p_0,p_1)=(p_0,-p_1)$ is the linear bijection \lean{ErdosSzekeres.Suk.reflectY}
(\dirn{ErdosSzekeres/Suk/Reflection.lean:34}); the simplification lemmas \lean{Suk.reflectY_x} and
\lean{Suk.reflectY_y} (\dirn{ErdosSzekeres/Suk/Reflection.lean:42} and \dirn{:43}) give its two
coordinates, and \lean{Suk.reflectY_reflectY} (\dirn{ErdosSzekeres/Suk/Reflection.lean:44}) says that
$\sigma(\sigma(p))=p$. \lean{ErdosSzekeres.Suk.generalPosition_reflectY}
(\dirn{ErdosSzekeres/Suk/Reflection.lean:52}) transports general position from $S$ to $\sigma(S)$ (the
forward direction of (c) for $g=\sigma$), and \lean{ErdosSzekeres.Suk.hasConvexNGon_of_reflectY}
(\dirn{ErdosSzekeres/Suk/Reflection.lean:64}) says that $\mathrm{HasConvexNGon}\,n\,\sigma(S)$ implies
$\mathrm{HasConvexNGon}\,n\,S$ (the reverse direction of (d) for $g=\sigma$).
\end{remark}

\begin{lemma}[Carath\'eodory's theorem in the plane; written step]\label{fd:lem:caratheodory}
Let $Y\subseteq\R^2$ and $x\in\conv Y$. Then there is a finite, affinely independent set $Y'\subseteq Y$
with $\card{Y'}\le3$ and $x\in\conv Y'$.
\end{lemma}

\begin{proof}
By Lemma~\ref{fd:lem:hull}(b), $x$ is a convex combination of finitely many points of $Y$; if $T_0$ is the
finite set of these points, then $x\in\conv T_0$ by Lemma~\ref{fd:lem:hull}(c). Among all finite sets
$T\subseteq Y$ with $x\in\conv T$ choose one, $T=\{y_1,\dots,y_k\}$ with distinct $y_i$, of least
cardinality $k$. Then $k\ge1$ because $\conv\emptyset=\emptyset$. By Lemma~\ref{fd:lem:hull}(c),
$x=\sum_i\lambda_iy_i$ with $\lambda_i\ge0$ and $\sum_i\lambda_i=1$. Every $\lambda_j$ is positive: if
$k=1$ then $\lambda_1=1$, and if $k\ge2$ and $\lambda_j=0$, then $x$ is a convex combination of the points
of $T\setminus\{y_j\}$, so $x\in\conv(T\setminus\{y_j\})$, contradicting minimality.

$T$ is affinely independent. Otherwise there are reals $\mu_1,\dots,\mu_k$, not all zero, with
$\sum_i\mu_i=0$ and $\sum_i\mu_iy_i=0$. Some $\mu_i$ is positive (they sum to $0$ and are not all zero).
Let $t=\min\{\lambda_i/\mu_i:\mu_i>0\}>0$, attained at $i=j$, and put $\lambda'_i=\lambda_i-t\mu_i$. Then
$\lambda'_i\ge0$ for all $i$ (if $\mu_i\le0$ then $\lambda'_i\ge\lambda_i$; if $\mu_i>0$ then
$t\le\lambda_i/\mu_i$), $\lambda'_j=0$, $\sum_i\lambda'_i=1-t\cdot0=1$ and
$\sum_i\lambda'_iy_i=x-t\cdot0=x$. So $x$ is a convex combination of the points of $T\setminus\{y_j\}$,
contradicting minimality again (note $k\ge2$ here, since $\sum\lambda'_i=1$ forces some
$\lambda'_i\ne0$, $i\ne j$).

Finally $k\le3$: if $k\ge4$, the $k-1\ge3$ vectors $y_2-y_1,\dots,y_k-y_1$ in the two-dimensional space
$\R^2$ are linearly dependent, $\sum_{i\ge2}\nu_i(y_i-y_1)=0$ with $\nu\ne0$; then $\mu_1=-\sum_{i\ge2}\nu_i$
and $\mu_i=\nu_i$ ($i\ge2$) satisfy $\sum_i\mu_i=0$ and $\sum_i\mu_iy_i=\sum_{i\ge2}\nu_i(y_i-y_1)=0$, with
$\mu\ne0$, contradicting affine independence. Take $Y'=T$.

The Lean code uses Mathlib's version of this theorem (module \lean{Mathlib.Analysis.Convex.Caratheodory},
with \lean{Caratheodory.minCardFinsetOfMemConvexHull}), whose proof is the same minimal-cardinality
argument; the bound $\card{Y'}\le3$ is derived in Lean from affine independence as
$\card{Y'}\le\dim+1=3$.
\end{proof}

\begin{lemma}[Four-point criteria; (a) \leanref{ErdosSzekeres.Suk.convexPosition_of_subsets_card_le_four}{ErdosSzekeres/Suk/FourPoint.lean:10}; (b) \leanref{ErdosSzekeres.Suk.convexPosition_of_four_subsets}{ErdosSzekeres/Suk/FourPoint.lean:40}]\label{fd:lem:four-subsets}\raggedright
Let $s$ be a finite set of points.
\begin{enumerate}[label=(\alph*)]
\item If every $t\subseteq s$ with $\card t\le4$ is in convex position, then $s$ is in convex position.
\item If $4\le\card s$ and every $t\subseteq s$ with $\card t=4$ is in convex position, then $s$ is in
  convex position.
\end{enumerate}
No general-position assumption is made.
\end{lemma}

\begin{proof}
We follow the Lean proofs. (a) By Lemma~\ref{fd:lem:cp-hull} we must show $p\notin\conv(s\setminus\{p\})$
for every $p\in s$. Suppose $p\in\conv(s\setminus\{p\})$. By Lemma~\ref{fd:lem:caratheodory} there is a
finite $Y'\subseteq s\setminus\{p\}$ with $\card{Y'}\le3$ and $p\in\conv Y'$. Let $u=Y'\cup\{p\}\subseteq s$.
Since $p\notin Y'$, $\card u=\card{Y'}+1\le4$ and $u\setminus\{p\}=Y'$. By hypothesis $u$ is in convex
position, so $p\notin\conv(u\setminus\{p\})=\conv Y'$ by Lemma~\ref{fd:lem:cp-hull}, a contradiction.

(b) Let $t\subseteq s$ with $\card t\le4$. Since $\card s\ge4$, we can add $4-\card t$ points of
$s\setminus t$ to $t$ (there are $\card s-\card t\ge4-\card t$ of them) and obtain $u$ with
$t\subseteq u\subseteq s$ and $\card u=4$. By hypothesis $u$ is in convex position, hence so is $t$
(Lemma~\ref{fd:lem:mono}(a)). Now (a) applies.
\end{proof}

\begin{lemma}[Half-planes of the orientation; (a) \leanref{ErdosSzekeres.Suk.orient_affine_combination}{ErdosSzekeres/Suk/SupportingLines.lean:6}; (b) \leanref{ErdosSzekeres.Suk.convex_orient_positive}{ErdosSzekeres/Suk/SupportingLines.lean:11}; (c) \leanref{ErdosSzekeres.Suk.convex_orient_nonnegative}{ErdosSzekeres/Suk/SupportingLines.lean:21}]\label{fd:lem:halfplane}\raggedright
The orientation is affine in its third argument, and its half-planes are convex.
\begin{enumerate}[label=(\alph*)]
\item For $p,q,a,b\in\R^2$ and $u,v\in\R$ with $u+v=1$:
  $\orient(p,q,u\,a+v\,b)=u\,\orient(p,q,a)+v\,\orient(p,q,b)$.
\item For $p,q\in\R^2$ the set $\{r\in\R^2:0<\orient(p,q,r)\}$ is convex.
\item For $p,q\in\R^2$ the set $\{r\in\R^2:0\le\orient(p,q,r)\}$ is convex.
\end{enumerate}
\end{lemma}

\begin{proof}
(a) $\orient(p,q,r)=(q_0-p_0)(r_1-p_1)-(q_1-p_1)(r_0-p_0)$ is a linear function of $r-p$. For
$r=ua+vb$ with $u+v=1$ we have $r-p=u(a-p)+v(b-p)$, and the claim follows. (Lean closes the same
polynomial identity with the certificate $(p_1q_0-p_0q_1)(u+v-1)$.)

(b) Let $\orient(p,q,a)>0$, $\orient(p,q,b)>0$ and $u,v\ge0$ with $u+v=1$. By (a) the value at $ua+vb$ is
$u\,\orient(p,q,a)+v\,\orient(p,q,b)$. If $u=0$ then $v=1$ and this is $\orient(p,q,b)>0$; otherwise
$u>0$, the first term is positive and the second is nonnegative.

(c) By (a) the value at $ua+vb$ is a sum of two nonnegative terms.
\end{proof}

\begin{lemma}[Strict supporting lines; \leanref{ErdosSzekeres.Suk.convexPosition_of_strict_supports}{ErdosSzekeres/Suk/SupportingLines.lean:28}]\label{fd:lem:strict-supports}\sloppy
Let $s$ be a finite set of points such that for every $p\in s$ there is $q\in\R^2$ with
$0<\orient(p,q,r)$ for every $r\in s$ with $r\ne p$. Then $s$ is in convex position.
\end{lemma}

\begin{proof}
We follow the Lean proof. Let $p\in s$ and choose $q$ as in the hypothesis. The set
$H=\{r:0<\orient(p,q,r)\}$ is convex (Lemma~\ref{fd:lem:halfplane}(b)) and contains $s\setminus\{p\}$, so
$\conv(s\setminus\{p\})\subseteq H$ by Lemma~\ref{fd:lem:hull}(a). But $\orient(p,q,p)=0$
(Lemma~\ref{fd:lem:orient}(e)), so $p\notin H$ and hence $p\notin\conv(s\setminus\{p\})$. By
Lemma~\ref{fd:lem:cp-hull}, $s$ is in convex position.
\end{proof}

\begin{lemma}[A support identity; \leanref{ErdosSzekeres.Suk.orient_support_sum}{ErdosSzekeres/Suk/SupportingLines.lean:40}]\label{fd:lem:support-sum}\sloppy
For all $a,p,b,r\in\R^2$: $\orient(p,\,p+b-a,\,r)=\orient(a,p,r)+\orient(p,b,r)$.
\end{lemma}

\begin{proof}
For $w\in\R^2$, $\orient(p,p+w,r)=\det(w\mid r-p)$ is linear in $w$. With $w=b-a=(p-a)+(b-p)$ we get
$\orient(p,p+b-a,r)=\det(p-a\mid r-p)+\det(b-p\mid r-p)$. The second term is $\orient(p,b,r)$. For the
first, $\orient(a,p,r)=\det(p-a\mid r-a)=\det(p-a\mid (r-p)+(p-a))=\det(p-a\mid r-p)$, since a determinant
with two equal columns vanishes. (The Lean proof expands both sides.)
\end{proof}

\begin{lemma}[Four counterclockwise points; \leanref{ErdosSzekeres.Suk.convexPosition_four_of_positive}{ErdosSzekeres/Suk/SupportingLines.lean:47}]\label{fd:lem:four-positive}\sloppy
Let $a,b,c,d\in\R^2$ with
\[
  0<\orient(a,b,c),\qquad 0<\orient(a,b,d),\qquad 0<\orient(a,c,d),\qquad 0<\orient(b,c,d).
\]
Then the finite set $\{a,b,c,d\}$ is in convex position.
\end{lemma}

\begin{proof}
We follow the Lean proof, which exhibits a strict supporting line at each of the four points and applies
Lemma~\ref{fd:lem:strict-supports}. (The hypotheses imply, by Lemma~\ref{fd:lem:orient}(e), that the four
points are pairwise distinct, but the argument does not need this.) Write $o_{abc}=\orient(a,b,c)$,
$o_{abd}=\orient(a,b,d)$, $o_{acd}=\orient(a,c,d)$ and $o_{bcd}=\orient(b,c,d)$, all positive. By
Lemma~\ref{fd:lem:orient}(b) every cyclic rotation of the arguments leaves $\orient$ unchanged, and by
Lemma~\ref{fd:lem:orient}(e) $\orient$ vanishes when two arguments are equal. For $p\in\{a,b,c,d\}$ and
$r\in\{a,b,c,d\}$ with $r\ne p$, the point $r$ is one of the three letters other than $p$ (as a point, $r$
differs from $p$), so it suffices to check the three letters other than $p$. Each line below applies
Lemma~\ref{fd:lem:support-sum}.
\begin{itemize}
\item $p=a$, $q=a+b-d$: $\orient(a,q,r)=\orient(d,a,r)+\orient(a,b,r)$. For $r=b$:
  $\orient(d,a,b)+0=o_{abd}$. For $r=c$: $\orient(d,a,c)+o_{abc}=o_{acd}+o_{abc}$. For $r=d$:
  $0+o_{abd}$.
\item $p=b$, $q=b+c-a$: $\orient(b,q,r)=\orient(a,b,r)+\orient(b,c,r)$. For $r=a$:
  $0+\orient(b,c,a)=o_{abc}$. For $r=c$: $o_{abc}+0$. For $r=d$: $o_{abd}+o_{bcd}$.
\item $p=c$, $q=c+d-b$: $\orient(c,q,r)=\orient(b,c,r)+\orient(c,d,r)$. For $r=a$:
  $\orient(b,c,a)+\orient(c,d,a)=o_{abc}+o_{acd}$. For $r=b$: $0+\orient(c,d,b)=o_{bcd}$. For $r=d$:
  $o_{bcd}+0$.
\item $p=d$, $q=d+a-c$: $\orient(d,q,r)=\orient(c,d,r)+\orient(d,a,r)$. For $r=a$:
  $\orient(c,d,a)+0=o_{acd}$. For $r=b$: $\orient(c,d,b)+\orient(d,a,b)=o_{bcd}+o_{abd}$. For $r=c$:
  $0+\orient(d,a,c)=o_{acd}$.
\end{itemize}
All twelve values are positive, so each of $a,b,c,d$ has a strict supporting line, and
Lemma~\ref{fd:lem:strict-supports} shows that $\{a,b,c,d\}$ is in convex position.
\end{proof}

\begin{lemma}[Positive triples; \leanref{ErdosSzekeres.Suk.convexPosition_of_positive_triples}{ErdosSzekeres/Suk/OrientationClique.lean:9}]\label{fd:lem:pos-triples}\sloppy
Let $\iota$ be a type with a linear order $<$, let $p\colon\iota\to\R^2$ be injective, and let $s$ be a
finite subset of $\iota$ with $4\le\card s$. Suppose that $0<\orient(p(i),p(j),p(k))$ for all
$i,j,k\in s$ with $i<j$ and $j<k$. Then the finite set $p(s)=\{p(i):i\in s\}$ is in convex position.
\end{lemma}

\begin{proof}
We follow the Lean proof. Since $p$ is injective, $\card{p(s)}=\card s\ge4$, so by
Lemma~\ref{fd:lem:four-subsets}(b) it suffices to show that every $t\subseteq p(s)$ with $\card t=4$ is in
convex position. Let $u=\{i\in s:p(i)\in t\}$. Since $t\subseteq p(s)$, $p(u)=t$, and $\card u=\card t=4$ by
injectivity. List the elements of the finite subset $u$ of the linearly ordered set $\iota$ in increasing
order, $e_0<e_1<e_2<e_3$. Then $t=\{p(e_0),p(e_1),p(e_2),p(e_3)\}$, and the hypothesis applied to the
increasing triples $(e_0,e_1,e_2)$, $(e_0,e_1,e_3)$, $(e_0,e_2,e_3)$ and $(e_1,e_2,e_3)$ of elements of
$s$ gives exactly the four hypotheses of Lemma~\ref{fd:lem:four-positive} for
$(a,b,c,d)=(p(e_0),p(e_1),p(e_2),p(e_3))$. Hence $t$ is in convex position.
\end{proof}

\begin{remark}\label{fd:rem:pos-triples}
(i) Only $s$ is assumed finite; $\iota$ may be infinite, and injectivity is assumed for $p$ on all of
$\iota$. (ii) The Lean proof discharges the order facts $0<1<2<3$ in $\mathrm{Fin}\,4$ and the equality
of the set of all elements of $\mathrm{Fin}\,4$ with $\{0,1,2,3\}$ by \texttt{decide}, which the kernel
evaluates without any additional axiom. These are side facts inside the proof of a statement that is not
itself a finite computation, and the written proof above needs no machine computation; so the statement
is not computer-assisted in the sense of the rigor labels. (iii) The
hypothesis $4\le\card s$ is needed by the Lean proof (through Lemma~\ref{fd:lem:four-subsets}(b)). As a
variant, not part of the Lean statement: the conclusion also holds when $\card s\le3$, by
Lemma~\ref{fd:lem:small}(b) when $\card{p(s)}\le2$, and when $\card s=3$ because the three points then
have nonzero orientation, hence are pairwise distinct and not collinear
(Lemmas~\ref{fd:lem:orient}(e) and~\ref{fd:lem:collinear}(a)), so Lemma~\ref{fd:lem:small}(c) applies.
\end{remark}

\begin{lemma}[Negative triples; \leanref{ErdosSzekeres.Suk.convexPosition_of_negative_triples}{ErdosSzekeres/Suk/OrientationClique.lean:33}]\label{fd:lem:neg-triples}\sloppy
Let $\iota$ be a type with a linear order $<$, let $p\colon\iota\to\R^2$ be injective, and let $s$ be a
finite subset of $\iota$ with $4\le\card s$. Suppose that $\orient(p(i),p(j),p(k))<0$ for all
$i,j,k\in s$ with $i<j$ and $j<k$. Then $p(s)$ is in convex position.
\end{lemma}

\begin{proof}
We follow the Lean proof, which applies Lemma~\ref{fd:lem:pos-triples} to the reversed order. Let
$i\prec j$ mean $j<i$; this is a linear order on $\iota$. If $i,j,k\in s$ with $i\prec j\prec k$, then
$k<j<i$, so $\orient(p(k),p(j),p(i))<0$ by hypothesis, and by Lemma~\ref{fd:lem:orient}(d) (the
transposition of the first and third arguments has sign $-1$)
$\orient(p(i),p(j),p(k))=-\orient(p(k),p(j),p(i))>0$. Lemma~\ref{fd:lem:pos-triples} for $(\iota,\prec)$
gives the claim.
\end{proof}

\begin{lemma}[Positive cycles; \leanref{ErdosSzekeres.hasConvexNGon_of_positive_cycle}{ErdosSzekeres/ConvexCycleCuts.lean:11}]\label{fd:lem:pos-cycle}\sloppy
Let $N,k\in\N$ with $4\le k$. Let $p\colon\mathrm{Fin}\,N\to\R^2$ be injective, let
$e\colon\mathrm{Fin}\,k\to\mathrm{Fin}\,N$ be an embedding (an injective map; Lean type
\texttt{Fin k}~$\hookrightarrow$~\texttt{Fin N}), and let $P$ be a finite set of points with $p(e(i))\in P$
for every $i\in\mathrm{Fin}\,k$. Suppose that
\[
  0<\orient\bigl(p(e(i)),\,p(e(j)),\,p(e(l))\bigr)\qquad\text{for all } i,j,l\in\mathrm{Fin}\,k
  \text{ with } i<j \text{ and } j<l .
\]
Then $\mathrm{HasConvexNGon}\,k\,P$.
\end{lemma}

\begin{proof}
We follow the Lean proof. The map $q=p\circ e\colon\mathrm{Fin}\,k\to\R^2$ is injective, as a composition of
injective maps. Let $T=q(\mathrm{Fin}\,k)=\{q(0),\dots,q(k-1)\}$. Then $T\subseteq P$ by hypothesis and
$\card T=k$ by injectivity. Lemma~\ref{fd:lem:pos-triples}, applied to $\iota=\mathrm{Fin}\,k$ with its
natural order, the injective map $q$ and $s=\mathrm{Fin}\,k$ (so $\card s=k\ge4$), shows that $T=q(s)$ is
in convex position; its hypothesis is exactly the assumed positivity. So $T$ witnesses
$\mathrm{HasConvexNGon}\,k\,P$.
\end{proof}

\begin{remark}\label{fd:rem:pos-cycle}\sloppy
The hypothesis concerns all increasing triples of the ordering $e$, not only consecutive ones. The name
refers to the clause family \lean{ConvexCycleCut} (\dirn{ErdosSzekeres/ConvexCycleCuts.lean:8}), which is
used as cuts in the SMT coordinate-repair searches described in \dirn{docs/SAT.md:168}--\dirn{233}:
\lean{convexCycleCut_of_avoids}
(\dirn{ErdosSzekeres/ConvexCycleCuts.lean:31}) is the
contrapositive of Lemma~\ref{fd:lem:pos-cycle}: under the same hypotheses on $k$, $p$, $e$ and $P$
(without the positivity hypothesis), if $P$ is $k$-free then some $i<j<l$ in $\mathrm{Fin}\,k$ has
$\orient(p(e(i)),p(e(j)),p(e(l)))\le0$.
\end{remark}

\begin{lemma}[Separation of a point from a finite set; written step]\label{fd:lem:separation}
Let $F\subseteq\R^2$ be finite and $x\in\R^2$. Then $x\notin\conv F$ if and only if there is a linear
functional $\ell\colon\R^2\to\R$ (that is, $\ell(z)=w_0z_0+w_1z_1$ for some $w\in\R^2$) with
$\ell(x)>\ell(y)$ for every $y\in F$.
\end{lemma}

\begin{proof}
Write $\langle w,z\rangle=w_0z_0+w_1z_1$ and $\lVert z\rVert^2=\langle z,z\rangle$ (the Euclidean norm,
written out).

($\Leftarrow$, elementary.) Let $\ell(z)=\langle w,z\rangle$ satisfy $\ell(x)>\ell(y)$ for all $y\in F$.
The set $H=\{z:\ell(z)<\ell(x)\}$ is convex by Lemma~\ref{fd:lem:hull}(d) and contains $F$, so
$\conv F\subseteq H$ by Lemma~\ref{fd:lem:hull}(a); since $x\notin H$, $x\notin\conv F$. (If $F=\emptyset$
this is also clear from $\conv\emptyset=\emptyset$.)

($\Rightarrow$.) Suppose $x\notin\conv F$. If $F=\emptyset$, the functional $\ell=0$ satisfies the
condition vacuously. Otherwise write $F=\{y_1,\dots,y_m\}$ with $m\ge1$ distinct points and let
$\Delta=\{\lambda\in\R^m:\lambda_i\ge0,\ \sum_i\lambda_i=1\}$. The set $\Delta$ is closed and bounded, hence
compact by the Heine--Borel theorem (a closed and bounded subset of $\R^m$ is compact; W.~Rudin,
\emph{Principles of Mathematical Analysis}, 3rd ed., Theorem~2.41), and the function
$h(\lambda)=\lVert x-\sum_i\lambda_iy_i\rVert^2$ is a polynomial in $\lambda$, hence continuous; so $h$
attains its minimum on the nonempty compact set $\Delta$ at some $\lambda^*$ (a continuous real function on
a nonempty compact metric space attains its minimum; Rudin, op.\ cit., Theorem~4.16). Put
$q=\sum_i\lambda^*_iy_i$. By Lemma~\ref{fd:lem:hull}(c), $\conv F=\{\sum_i\lambda_iy_i:\lambda\in\Delta\}$,
so $q\in\conv F$ and
\begin{equation*}
  \lVert x-q\rVert^2\le\lVert x-z\rVert^2\qquad\text{for every }z\in\conv F. \tag{$**$}
\end{equation*}
Since $x\notin\conv F$, $x\ne q$; put $w=x-q\ne0$ and $\ell(z)=\langle w,z\rangle$. Let $y\in F$ and
$0<t\le1$. The point $z_t=(1-t)q+ty$ lies in $\conv F$ (convexity, Lemma~\ref{fd:lem:hull}(a)), and
$x-z_t=w-t(y-q)$, so by $(**)$
\[
  \lVert w\rVert^2\le\lVert w-t(y-q)\rVert^2=\lVert w\rVert^2-2t\langle w,y-q\rangle+t^2\lVert y-q\rVert^2 .
\]
Hence $2\langle w,y-q\rangle\le t\lVert y-q\rVert^2$ for every $t\in(0,1]$, and letting $t\to0$ gives
$\langle w,y-q\rangle\le0$, that is, $\ell(y)\le\ell(q)$. Finally
$\ell(x)-\ell(q)=\langle w,x-q\rangle=\lVert w\rVert^2>0$, so $\ell(y)\le\ell(q)<\ell(x)$.
\end{proof}

\paragraph{Coverage.}
{\raggedright
\begin{itemize}
\item \emph{Covered definitions, with the label of the statement.} \lean{ErdosSzekeres.Point}
  (\ref{fd:def:point}), \lean{ErdosSzekeres.orient} (\ref{fd:def:orient}),
  \lean{ErdosSzekeres.GeneralPosition} (\ref{fd:def:gp}), \lean{ErdosSzekeres.ConvexPosition}
  (\ref{fd:def:cp}), \lean{ErdosSzekeres.HasConvexNGon} (\ref{fd:def:ngon}),
  \lean{ErdosSzekeres.ForcesConvexNGon} (\ref{fd:def:forces}), the three copies of \lean{Invariant}
  (\ref{fd:def:inv}) and the three cited copies of \lean{rotAbout} (\ref{fd:def:rot}), together with the
  Mathlib notions \lean{Convex}, \lean{convexHull}, \lean{AffineIndependent}, \lean{AffineMap},
  \lean{AffineEquiv} (\ref{fd:def:conv}), \lean{Collinear} (\ref{fd:def:gp}) and \lean{ConvexIndependent}
  (\ref{fd:def:cp}).
\item \emph{Covered facts, with the label of the statement.} \lean{orient_swap}, \lean{orient_cycle},
  \lean{orient_translate} (\ref{fd:lem:orient}); \lean{SymReduction.orient_affine} and, for part (i),
  \lean{OrbitConsistency.orient_rotAbout_inv} (\ref{fd:lem:affine-orient});
  \lean{orient_eq_zero_iff_collinear}, \lean{generalPosition_iff_orient_ne_zero} (\ref{fd:lem:collinear});
  \lean{convexPosition_iff} (\ref{fd:lem:cp-hull}); \lean{ConvexPosition.mono}, \lean{GeneralPosition.mono},
  \lean{HasConvexNGon.mono}, \lean{MirrorOdd.hasConvexNGon_of_le} (\ref{fd:lem:mono});
  \lean{HasConvexNGon.card_le}, \lean{ConvexPosition.hasConvexNGon} (\ref{fd:lem:ngon-basic});
  \lean{ForcesConvexNGon.mono}, \lean{forcesConvexNGon_iff_exact} (\ref{fd:lem:forces-basic});
  \lean{convexPosition_of_card_three} (\ref{fd:lem:small}); \lean{SymReduction.convexPosition_image_affine},
  \lean{SymReduction.generalPosition_image_affine}, \lean{SymReduction.generalPosition_image_affineEquiv},
  \lean{SymReduction.hasConvexNGon_image_affine}, \lean{SymReduction.injective_of_affine}
  (\ref{fd:lem:affine-cp}, parts (b$'$)--(e$'$)); \lean{Suk.convexPosition_of_subsets_card_le_four},
  \lean{Suk.convexPosition_of_four_subsets} (\ref{fd:lem:four-subsets}); \lean{Suk.orient_affine_combination},
  \lean{Suk.convex_orient_positive}, \lean{Suk.convex_orient_nonnegative} (\ref{fd:lem:halfplane});
  \lean{Suk.convexPosition_of_strict_supports} (\ref{fd:lem:strict-supports}); \lean{Suk.orient_support_sum}
  (\ref{fd:lem:support-sum}); \lean{Suk.convexPosition_four_of_positive} (\ref{fd:lem:four-positive});
  \lean{Suk.convexPosition_of_positive_triples} (\ref{fd:lem:pos-triples});
  \lean{Suk.convexPosition_of_negative_triples} (\ref{fd:lem:neg-triples});
  \lean{hasConvexNGon_of_positive_cycle} (\ref{fd:lem:pos-cycle}); \lean{LayerLemma.iterate_mem}
  (\ref{fd:lem:inv-image}); \lean{LayerLemma.rotAbout_eq_crot}, \lean{LayerLemma.rotAff},
  \lean{LayerLemma.rotAbout_iterate}, \lean{LayerLemma.rotAbout_injective} (\ref{fd:lem:rot}).
\item \emph{Written steps.} The following statements and parts carry no Lean tag in this section, and
  their proofs above are written steps. We found no declaration in the files of the table of sources that
  states one of them in the form given here. Several of them have Lean counterparts in another form,
  special cases or variants, in the project or in Mathlib; we name some of these, here or in the remark
  indicated, but the research records contain many further copies and variants of such elementary facts,
  and we do not list them all. Naming a counterpart does not make a written step Lean-checked in the sense
  of the rigor labels; we give axioms evidence for such a counterpart only where we say so. The written
  steps are:
  \begin{itemize}
  \item \ref{fd:lem:orient}(d) for a general commutative ring (its real case is
    \lean{OrbitConsistency.orient3_perm}, Remark~\ref{fd:rem:orient-lean}) and (e);
  \item part (ii) of \ref{fd:lem:affine-orient} (its case $o=0$, $\varphi=0$ is \lean{Suk.orient_reflectY},
    Remark~\ref{fd:rem:affine-orient-lean}; with the line given by a unit direction vector instead of the
    angle $\varphi$, (ii) is stated in the research file \dirn{H2/D3Encoding/Geometry.lean} as
    \lean{D3.orient_reflC} (\dirn{:87}));
  \item \ref{fd:lem:collinear-mem} (Mathlib's \lean{collinear_iff_of_mem}) and \ref{fd:lem:hull};
  \item \ref{fd:lem:inv-image}(a),(b) (the argument of (a) is a step of the proof of
    \lean{LayerLemma.map_mem_hullVerts});
  \item in \ref{fd:lem:rot}, the matrix form of Definition~\ref{fd:def:rot} and the identity
    $\rho_\theta(c)=c$ in (b) (the latter a step of the proof of \lean{LayerLemma.invariant_erase_centre}),
    the composition law and $\rho_0=\mathrm{id}$ in (c), the bijectivity and the inverse in (d), and (e),
    (f) (the identity $\rho_\theta(c)=c$ and the identity of (e) are stated, for the copy of
    \lean{rotAbout} in the review file \dirn{SYMLEAN/review/ReviewSanity.lean}, as
    \lean{SymReduction.Review.rotAbout_center} (\dirn{:1086}) and \lean{SymReduction.Review.orient_rotAbout}
    (\dirn{:1091}, which writes $(p_0-c_0)^2+(p_1-c_1)^2$ as \texttt{eDist2 p c}); we do not cite them as
    Lean-checked, because none of the theorems whose axioms the saved output
    \dirn{SYMLEAN/logs/ReviewSanity.out} of that file prints uses them);
  \item \ref{fd:lem:small}(a)--(c) apart from its last assertion;
  \item in \ref{fd:lem:affine-cp}, part (a) (Mathlib's \lean{AffineMap.image_convexHull}), the implication
    ``injective $\Rightarrow$ $\det A\ne0$'' with the inverse, and the reverse directions of (b)--(d) (for
    the reflection $(p_0,p_1)\mapsto(p_0,-p_1)$, the reverse direction of (d) is
    \lean{Suk.hasConvexNGon_of_reflectY}; Remark~\ref{fd:rem:affine-cp-lean} also names the
    forward-direction special cases \lean{HasConvexNGon.linearEquiv} and \lean{Suk.generalPosition_reflectY});
  \item \ref{fd:lem:caratheodory} (Mathlib's Carath\'eodory module);
  \item \ref{fd:lem:separation} (a variant of its easy direction is the main-library theorem
    \lean{Omission.BoundaryGeometry.exposed_not_in_hull},
    \dirn{ErdosSzekeres/Omission/BoundaryGeometry.lean:30}: if $S$ is a set of points, $p$ a point and
    $A,B,C$ reals with $Ap_0+Bp_1+C=0$ and $Aq_0+Bq_1+C>0$ for every $q\in S\setminus\{p\}$, then
    $p\notin\conv(S\setminus\{p\})$).
  \end{itemize}
\item \emph{Other declarations named.} Other declarations of the project that are named in statements,
  remarks, proofs or this paragraph as Lean counterparts, as definitions and lemmas that these rely on, or
  as the location of a proof step: \lean{OrbitConsistency.orient3_perm} and \lean{OrbitConsistency.orient3},
  \lean{Suk.orient_reflectY}, \lean{HasConvexNGon.linearEquiv}, \lean{Suk.reflectY}, \lean{Suk.reflectY_x},
  \lean{Suk.reflectY_y}, \lean{Suk.reflectY_reflectY}, \lean{Suk.generalPosition_reflectY},
  \lean{Suk.hasConvexNGon_of_reflectY}, \lean{ConvexCycleCut}, \lean{convexCycleCut_of_avoids},
  \lean{SymReduction.affineMap_eq_mulVec}, \lean{LayerLemma.map_mem_hullVerts},
  \lean{LayerLemma.invariant_erase_centre}, \lean{D3.orient_reflC}, \lean{SymReduction.Review.rotAbout_center},
  \lean{SymReduction.Review.orient_rotAbout} and \lean{Omission.BoundaryGeometry.exposed_not_in_hull}.
\item \emph{Deliberately not covered.} Of the declarations of the main-library files in the table of
  sources, all others are listed here: in \dirn{ErdosSzekeres/Definitions.lean}, the bound function
  \lean{conjecturedBound} ($n\mapsto2^{\max(n-2,0)}+1$) and the open conjectures
  \lean{ErdosSzekeresConjecture} and \lean{HeptagonConjecture} (propositions that are defined but not
  asserted), none of which is used by the facts of this section; in \dirn{ErdosSzekeres/Orientation.lean}
  and \dirn{ErdosSzekeres/OrientationGeometry.lean}, \lean{rationalToReal}, \lean{orient_rationalToReal},
  \lean{orient_rationalToReal_ne_zero}, \lean{RationalGeneralPosition} and
  \lean{generalPosition_of_rational} (transfer of exact rational certificates to real points, outside the
  foundations list), and the private helper \lean{exists_smul_of_det_zero} (absorbed into the proof of
  Lemma~\ref{fd:lem:collinear}); in \dirn{ErdosSzekeres/Basic.lean},
  \lean{convexIndependent_of_affineIndependent} (a bridge to Mathlib's affine independence, used in the
  main library only for \lean{convexPosition_of_card_three}, which we prove directly) and
  \lean{forces_triangle}, \lean{conjecture_case_three}, \lean{conjecture_implies_seven},
  \lean{counterexample_refutes_conjecture} (small corollaries about the conjecture, not used by the facts
  of this section; \lean{forces_triangle} is named above only in the axioms evidence); and in
  \dirn{ErdosSzekeres/Suk/OrientationClique.lean}, \lean{convexPosition_insert_of_positive_triples} and
  \lean{convexPosition_insert_of_negative_triples} (outside the foundations list). Of the research files
  in the table, only the declarations named in the statements, remarks or proofs of this section belong to
  its scope; the theorems named in the axioms-evidence paragraph only as evidence, and the other
  declarations of these files, are treated in other sections or lie outside Stage~1. In particular,
  \lean{toC}, \lean{ofC}, \lean{crot} and their lemmas in \dirn{LAYERLEAN/LayerLemma.lean} and its copies
  are implementation devices for complex coordinates, replaced here by the direct complex form in
  Lemma~\ref{fd:lem:rot}(a).
\end{itemize}
\par}
